% !TEX root = ../../main.tex
\setchapterimage{}
\setchapterintro{从环同态和理想的对应出发，讨论环的直和、分式环与多项式环，并研究整环及域上多项式的性质。}
\chapter{环}
\label{p2:ch3}

\section{环的同态定理}

\begin{proposition}
若$H,N$为环$R$的子环，则$H\cap N$为子环；

若$N$为$R$的理想，则$H\cap N$为$H$的理想。

更进一步，若$N,H$都是$R$的理想，那么$H\cap N$为$R$的理想。


\end{proposition}
\begin{proof}
若$x,y\in H\cap N$，则$x-y,xy$同时属于$H$和$N$，故$H\cap N$为$R$的子环。

若$N$为$R$的理想，则对$x\in H\cap N$、$h\in H$，有$hx,xh\in N$；又因$H$为子环，$hx,xh\in H$。因此
\[
hx,xh\in H\cap N,
\]
即$H\cap N$为$H$的理想。

若$H,N$均为$R$的理想，则对任意$r\in R$和$x\in H\cap N$，$rx,xr$同时属于$H$与$N$，所以
\[
rx,xr\in H\cap N.
\]
故$H\cap N$为$R$的理想。


\end{proof}

\begin{definition}
	定义
\[
H+N=\{a+b:a\in H,\ b\in N\}.
\]
\end{definition}
\begin{proposition}
	$H+N $是$R$的加法子群。若$N$为$R$的理想，则$H+N$为$R$的子环；若$H,N$均为$R$的理想，则$H+N$也是理想。
\end{proposition}
\begin{proof}
	集合$H+N$对加法和取负封闭。若$N$为理想，取$a_i\in H,b_i\in N$，则
	\begin{align*}
		(a_1+b_1)(a_2+b_2)
		&=a_1a_2+a_1b_2+b_1a_2+b_1b_2\in H+N.
	\end{align*}
	故$H+N$为子环。若$H$也为理想，则对$r\in R$有
	\[
	r(a+b)=ra+rb\in H+N,\qquad (a+b)r=ar+br\in H+N,
	\]
	所以$H+N$为理想。
\end{proof}
\begin{example}
令$R=M_2(F)$，
\[
H=\left\{\begin{pmatrix}0&0\\a&0\end{pmatrix}:a\in F\right\},\qquad
N=\left\{\begin{pmatrix}0&b\\0&0\end{pmatrix}:b\in F\right\}.
\]
$H,N$都是子环，但$H+N$不是子环：取其中两个非零矩阵相乘，可得到非零对角矩阵。
\end{example}

\begin{lemma}\label{p2:ring-image-preimage}
设$\sigma:R\to R'$为环同态，$N=\ker\sigma$，$H$为$R$的子环，$H'$为$R'$的子环。则
\begin{align*}
\sigma^{-1}(\sigma(H))&=H+N,\\
\sigma(\sigma^{-1}(H'))&=H'\cap\sigma(R).
\end{align*}
特别地，第一式等于$H$当且仅当$N\subseteq H$；若$\sigma$满射，则第二式等于$H'$。
\end{lemma}
\begin{proof}
\textbf{第一式：}先证$\sigma^{-1}(\sigma(H))\subseteq H+N$。任取$x\in\sigma^{-1}(\sigma(H))$，则$\sigma(x)\in\sigma(H)$，所以存在$h\in H$使$\sigma(x)=\sigma(h)$。于是
\[
\sigma(x-h)=\sigma(x)-\sigma(h)=0,
\]
即$x-h\in N$。因此
\[
x=h+(x-h)\in H+N.
\]

再证$H+N\subseteq\sigma^{-1}(\sigma(H))$。任取$x\in H+N$，写作$x=h+n$，其中$h\in H$、$n\in N$。由于$\sigma(n)=0$，
\[
\sigma(x)=\sigma(h)+\sigma(n)=\sigma(h)\in\sigma(H),
\]
故$x\in\sigma^{-1}(\sigma(H))$。合并两个包含关系，得到第一式。此式不要求$\sigma$满射。

因为$N$是理想，$H+N$为子环；任一同时包含$H$和$N$的子环都包含它们的和。因此$\sigma^{-1}(\sigma(H))$正是包含$H$和$\ker\sigma$的最小子环。又因$N\subseteq H+N$，
\[
\sigma^{-1}(\sigma(H))=H
\quad\Longleftrightarrow\quad H+N=H
\quad\Longleftrightarrow\quad N\subseteq H.
\]

\textbf{第二式：}任取$y\in\sigma(\sigma^{-1}(H'))$，则存在$x\in\sigma^{-1}(H')$使$y=\sigma(x)$。由反像的定义，$\sigma(x)\in H'$，故
\[
y\in H'\cap\sigma(R).
\]
反过来，若$y\in H'\cap\sigma(R)$，则存在$x\in R$使$y=\sigma(x)$；又因$y\in H'$，有$x\in\sigma^{-1}(H')$，所以
\[
y=\sigma(x)\in\sigma(\sigma^{-1}(H')).
\]
这就证明第二式。若$\sigma$满射，$\sigma(R)=R'$，于是
\[
\sigma(\sigma^{-1}(H'))=H'\cap R'=H'.
\]
\end{proof}

\begin{theorem}\label{p2:ring-correspondence}
\textbf{环的对应定理：}设$\sigma:R\to R'$为满环同态，$N=\ker\sigma$。映射$H\mapsto\sigma(H)$给出$R$中包含$N$的子环与$R'$的子环之间的一一对应；在此对应下，理想与理想对应。
\end{theorem}
\begin{proof}
\textbf{$\sigma(H)$是子环：}若$H$为子环，则$\sigma(H)$对减法和乘法封闭。

\textbf{$H$是理想，则$\sigma(H)$是理想：}若$H$为理想，任取$r'\in R'$及$a'\in\sigma(H)$，由满射取$r\in R,a\in H$使$\sigma(r)=r'$、$\sigma(a)=a'$，于是
\[
r'a'=\sigma(ra)\in\sigma(H),\qquad a'r'=\sigma(ar)\in\sigma(H).
\]
故$\sigma(H)$为理想。

\textbf{子环$H'\subseteq R'$的完全反像$\sigma^{-1}(H')$为子环：}反过来，子环$H'\subseteq R'$的完全反像$\sigma^{-1}(H')$是$R$的一个加法子群且对乘法封闭，从而为子环且包含$N$。

\textbf{理想$H'\subseteq R'$的完全反像$\sigma^{-1}(H')$为理想：}若$H'$为理想，任取$r\in R$、$a\in\sigma^{-1}(H')$，则
\[
\sigma(ra)=\sigma(r)\sigma(a)\in H',\qquad
\sigma(ar)=\sigma(a)\sigma(r)\in H',
\]
故$\sigma^{-1}(H')$为理想。

由命题\ref{p2:ring-image-preimage}，当$N\subseteq H$且$\sigma$满射时，
\[
\sigma^{-1}(\sigma(H))=H+N=H,\qquad
\sigma(\sigma^{-1}(H'))=H'.
\]
因此$H\mapsto\sigma(H)$与$H'\mapsto\sigma^{-1}(H')$互逆，得到所述一一对应。
\end{proof}

\begin{theorem}\label{p2:ring-third-iso}
设$\sigma:R\to R'$为环的满同态，$N=\ker\sigma$，$H$为$R$的任一包含$N$的理想。则$\sigma$诱导出环同构
\[
\overline\sigma:R/H\longrightarrow R'/\sigma(H),
\qquad
\overline\sigma(x+H)=\sigma(x)+\sigma(H).
\]
特别地，若取$R'=R/N$，$\sigma(x)=x+N$，则$\sigma(H)=H/N$，从而
\[
R/H\cong(R/N)/(H/N).
\]
同构映射为
\[
x+H\longmapsto(x+N)+(H/N).
\]
\end{theorem}
\begin{proof}
记$H'=\sigma(H)$。由环的对应定理，$H'$为$R'$的理想。由加法群的同构定理，$\sigma$诱导出加法群同构
\[
\overline\sigma:R/H\longrightarrow R'/H',
\qquad x+H\longmapsto\sigma(x)+H'.
\]

再验证$\overline\sigma$保持乘法。任取$x,y\in R$，有
\begin{align*}
\overline\sigma\bigl((x+H)(y+H)\bigr)
&=\overline\sigma(xy+H)\\
&=\sigma(xy)+H'\\
&=\sigma(x)\sigma(y)+H'\\
&=(\sigma(x)+H')(\sigma(y)+H')\\
&=\overline\sigma(x+H)\,\overline\sigma(y+H).
\end{align*}
所以$\overline\sigma$是环同构。

取自然同态$\sigma:R\to R/N$，$\sigma(x)=x+N$，此时$\sigma(H)=H/N$。代入上述同构，即得特别情形及其同构映射。
\end{proof}

\begin{theorem}\label{p2:ring-second-iso}
若$H$为$R$的子环，$N$为$R$的理想，则
\[
H/(H\cap N)\cong(H+N)/N,\qquad
h+(H\cap N)\longmapsto h+N.
\]
\end{theorem}
\begin{proof}
将自然同态$\pi:R\to R/N$限制在$H$上。其核为$H\cap N$，故
\[
H/(H\cap N)\cong\pi(H).
\]
另一方面，$H+N$为子环且包含$N$，并有$\pi(H+N)=\pi(H)$。将$\pi$限制在$H+N$上，其核为$N$，于是
\[
(H+N)/N\cong\pi(H).
\]
合并两个同构得到结论；所列映射的逆映射由$h+n+N\mapsto h+(H\cap N)$给出。
\end{proof}
\section{环的直和}

本节的环均有单位元,也就是说默认为幺环.
\begin{definition}
设$R_1,\ldots,R_r$为环。它们的加法群直和$R_1\oplus\cdots\oplus R_r$上再定义
\[
(a_1,\ldots,a_r)(b_1,\ldots,b_r)=(a_1b_1,\ldots,a_rb_r).
\]
所得环称为$R_1,\ldots,R_r$的\textbf{直和},记为$R=\bigoplus\limits_{i=1}^r R_i$,零元素为$(0_{R_1},\ldots,0_{R_r})$,单位元素为$(1_{R_1},\ldots,1_{R_r})$。
若$R_1,\ldots,R_r$都是交换环，那么$R$也是交换环.
\end{definition}


\begin{proposition}
在直和$R=R_1\oplus\cdots\oplus R_r$中，令
\[
R_i'=\{(0,\ldots,0,a_i,0,\ldots,0):a_i\in R_i\}.
\]
则每个$R_i'$都是$R$的理想，且$R_i'\cong R_i$；各$R_i'$的和为$R$，其中每个元素的分解唯一，并且$i\ne j$时$R_i'R_j'=(0)$。
\end{proposition}
\begin{proof}[补充证明]
第$i$坐标的嵌入
\[
a_i\longmapsto(0,\ldots,0,a_i,0,\ldots,0)
\]
保持加法和乘法，且其像为$R_i'$，故$R_i'\cong R_i$。

对$y=(b_1,\ldots,b_r)\in R$和$x=(0,\ldots,a_i,\ldots,0)\in R_i'$，按分量相乘得到
\[
yx=(0,\ldots,b_ia_i,\ldots,0),\qquad
xy=(0,\ldots,a_ib_i,\ldots,0).
\]
两者都属于$R_i'$，所以$R_i'$为理想。每个$y$都唯一写作
\[
y=\sum_{i=1}^r(0,\ldots,0,b_i,0,\ldots,0).
\]
最后，当$i\ne j$时，$R_i'$与$R_j'$的元素逐分量相乘得到零元。
\end{proof}

\begin{theorem}\label{p2:ring-internal-sum}
设$R_1,\ldots,R_r$为$R$的子环，且每个$R_i$均为$R$的理想。若
\[
R=R_1+\cdots+R_r,\qquad
R_i\cap(R_1+\cdots+\widehat{R_i}+\cdots+R_r)=(0)
\quad(1\leq i\leq r),
\]
则乘法映射给出环同构
\[
R_1\oplus\cdots\oplus R_r\cong R.
\]
此时称$R$为这些子环的\textbf{内直和}。
\end{theorem}
\begin{proof}[补充证明]
由加法群直和定理，每个$a\in R$可唯一写成$a=a_1+\cdots+a_r$，其中$a_i\in R_i$。

对$i\ne j$，因$R_i,R_j$均为理想，$a_ib_j$同时属于$R_i$与$R_j$，故
\[
a_ib_j\in R_i\cap R_j=(0).
\]
于是
\[
(a_1+\cdots+a_r)(b_1+\cdots+b_r)
=a_1b_1+\cdots+a_rb_r.
\]
加法群的同构因此也保持乘法。
\end{proof}

\begin{definition}
若$H,N$为$R$的理想，记$HN$为所有有限和$\sum_i a_ib_i$（$a_i\in H,b_i\in N$）组成的集合，称为两理想的\textbf{积}。若$H+N=R$，称$H,N$\textbf{互素}。
\end{definition}
\begin{proof}[补充证明]
$HN$对加法及取负封闭。对$r\in R$，
\[
r\sum_i a_ib_i=\sum_i(ra_i)b_i\in HN,\qquad
\biggl(\sum_i a_ib_i\biggr)r=\sum_i a_i(b_ir)\in HN.
\]
所以$HN$为理想。乘法的分配律还给出
\[
H(N+K)=HN+HK,\qquad (N+K)H=NH+KH.
\]
\end{proof}

\begin{lemma}\label{p2:ring-comaximal-product}
设$H,N,K$为$R$的理想。
若$R$交换且$H+N=R$，则$HN=H\cap N$。
若$H+N=R$且$K+N=R$，则$HK+N=R$。
\end{lemma}
\begin{proof}
总有$HN\subseteq H\cap N$。在交换情形，取$a\in H,b\in N$使$a+b=1$。若$c\in H\cap N$，则
\[
c=c(a+b)=ca+cb\in HN,
\]
所以$H\cap N\subseteq HN$。

再取$a\in H,b\in K,c,d\in N$使$a+c=1$、$b+d=1$。相乘得到
\[
1=ab+ad+cb+cd\in HK+N.
\]
因此$HK+N=R$。
\end{proof}

\begin{theorem}\label{p2:ring-crt}
设$R$为交换幺环，理想$N_1,\ldots,N_r$两两互素。则
\[
R/(N_1\cap\cdots\cap N_r)\cong R/N_1\oplus\cdots\oplus R/N_r,
\]
同构由$x+(N_1\cap\cdots\cap N_r)\mapsto(x+N_1,\ldots,x+N_r)$给出。
\end{theorem}
\begin{proof}
当$r=1$时结论即$R/N_1\cong R/N_1$。以下设$r\geq2$。

令$\sigma:R\to\bigoplus_iR/N_i$为各自然同态的乘积，则
\[
\ker\sigma=N_1\cap\cdots\cap N_r.
\]
需证$\sigma$满射。记
\[
M_i=N_1\cdots N_{i-1}N_{i+1}\cdots N_r.
\]
由理想积对和的分配律及前一引理，
\begin{align*}
M_1+M_2
&=(N_2+N_1)N_3\cdots N_r=N_3\cdots N_r,\\
M_1+M_2+M_3
&=(N_3+N_1N_2)N_4\cdots N_r=N_4\cdots N_r.
\end{align*}
继续如此，得到$M_1+\cdots+M_r=R$。故可选$e_i\in M_i$使
\[
e_1+\cdots+e_r=1.
\]
当$j\ne i$时，$e_j\in N_i$；对上式模$N_i$取像，可知$e_i\equiv1\pmod{N_i}$。

给定$x_1,\ldots,x_r\in R$，令$x=\sum_i e_ix_i$。则对每个$i$，
\[
x\equiv x_i\pmod{N_i}.
\]
因此$\sigma$满射，结论由环同态基本定理得到。
\end{proof}

\begin{corollary}\label{p2:ring-crt-congruence}
在上一定理的条件下，同余方程组
\[
x\equiv b_i\pmod{N_i}\qquad(1\leq i\leq r)
\]
有解，且任意两解模$N_1\cap\cdots\cap N_r$同余。
若交为$(0)$，则$R\cong\bigoplus_i R/N_i$。
\end{corollary}
\begin{proof}
存在性就是映射$\sigma$的满射性。两个解$x,y$满足全部同余式，当且仅当$x-y$属于每个$N_i$，即属于它们的交。交为零时，$\sigma$还单射。
\end{proof}

\begin{example}
若$n=\prod_{i=1}^r p_i^{e_i}$为素数幂分解，则理想$(p_i^{e_i})$两两互素，且它们的交为$(n)$。于是
\[
\Z/(n)\cong\bigoplus_{i=1}^r\Z/(p_i^{e_i}).
\]
\end{example}

\begin{definition}
对非空子集$S\subseteq R$，所有包含$S$的理想之交称为由$S$\textbf{生成的理想}，记为$(S)$。若$S=\{a_1,\ldots,a_r\}$，简记为$(a_1,\ldots,a_r)$。可由有限多个元素生成的理想称为\textbf{有限生成理想}；由一个元素生成的理想称为\textbf{主理想}。
\end{definition}
\begin{proof}[补充证明]
对$a\in R$，由$a$生成的理想恰为所有
\[
na+xa+ay+\sum_{i=1}^m x_iay_i,
\qquad n\in\Z,\quad x,y,x_i,y_i\in R,
\]
组成的集合。这些元素对加法、取负及左右乘以$R$的元素均封闭，并包含$a$；而任意包含$a$的理想也包含所有这些项，故它就是$(a)$。

若$R$交换，上式化为$na+xa$。若$R$还有单位元素，则$na=(n1)a$，因而$(a)=Ra$。同样，
\[
(a_1,\ldots,a_r)=(a_1)+\cdots+(a_r).
\]
\end{proof}
%\section{环的反同构}
%
%\begin{definition}
%两个环$R,R'$之间的双射$\sigma$若满足
%\[
%\sigma(a+b)=\sigma(a)+\sigma(b),\qquad
%\sigma(ab)=\sigma(b)\sigma(a),
%\]
%则称为\textbf{反同构}；当$R'=R$时称为\textbf{反自同构}。
%不要求双射时，满足这两式的映射称为\textbf{反同态}。
%\end{definition}
%
%\begin{example}
%对域$F$上的方阵环$M_n(F)$，转置满足
%\[
%(A+B)^T=A^T+B^T,\qquad (AB)^T=B^TA^T.
%\]
%它是反自同构，且连续转置两次得到原矩阵。
%\end{example}
%
%\begin{example}
%四元数体$\mathbb H$中，若$\alpha=a_0+a_1I+a_2J+a_3K$，定义
%\[
%\overline\alpha=a_0-a_1I-a_2J-a_3K.
%\]
%由$I^2=J^2=K^2=-1$、$IJ=-JI=K$、$JK=-KJ=I$、$KI=-IK=J$，可直接检查
%\[
%\overline{\alpha\beta}=\overline\beta\,\overline\alpha,
%\qquad \overline{\overline\alpha}=\alpha.
%\]
%因此共轭是$\mathbb H$的反自同构。
%\end{example}
%\begin{proof}
%共轭保持加法，且由定义可知$\overline{\overline\alpha}=\alpha$，所以它是双射。
%
%按$1,I,J,K$的乘法表逐项验证乘法次序反转。例如在$\alpha\beta$中，$a_1b_2IJ=a_1b_2K$经共轭变为$-a_1b_2K$；在$\overline\beta\,\overline\alpha$中，对应项为
%\[
%(-b_2J)(-a_1I)=a_1b_2JI=-a_1b_2K.
%\]
%其余基底乘积同理。由乘法对加法的分配律，$\overline{\alpha\beta}=\overline\beta\,\overline\alpha$对任意$\alpha,\beta$成立。
%\end{proof}
%
%\begin{example}
%对$M_2(\mathbb H)$，先逐项取四元数共轭，再转置，所得映射$A\mapsto\overline A^{T}$也为反自同构。
%\end{example}
%\begin{proof}[补充证明]
%设$A=(a_{ij})$、$B=(b_{ij})$。其$(i,j)$项满足
%\begin{align*}
%\bigl(\overline{AB}^{T}\bigr)_{ij}
%&=\overline{(AB)_{ji}}
%=\overline{\sum_k a_{jk}b_{ki}}\\
%&=\sum_k\overline{b_{ki}}\,\overline{a_{jk}}
%=\bigl(\overline B^{T}\overline A^{T}\bigr)_{ij}.
%\end{align*}
%加法与双射性也逐项成立。
%\end{proof}
%
%\begin{proposition}
%环$R$的自同构和反自同构在复合下组成一个群，$\Aut(R)$在其中的指数至多为$2$。若$R$非交换且存在反自同构，则指数为$2$。
%\end{proposition}
%\begin{proof}[补充证明]
%同型映射复合仍同型，反同型映射与同型映射复合为反同型映射，两个反同型映射复合为同型映射。逆映射保持原来的类型，所以这些映射组成群。
%
%若存在反自同构$\tau$，则每个反自同构$\rho$都可写成$(\rho\tau^{-1})\tau$，其中$\rho\tau^{-1}$为自同构；故反自同构至多构成$\Aut(R)\tau$这一陪集，指数至多为$2$。
%
%若$R$非交换，反自同构不可能同时是自同构：否则满射性将给出$ab=ba$对所有$a,b\in R$成立。因此在这一情形下两个陪集不同，指数恰为$2$。
%\end{proof}
%
%\begin{definition}
%对任意环$R$，在同一个加法群上定义新乘法$a\star b=ba$。得到的环记为$R^{\mathrm{op}}$，称为$R$的\textbf{反环}。恒等集合映射$R\to R^{\mathrm{op}}$是反同构。
%\end{definition}
%
%\section{素理想和极大理想}
%
%从本节起，除另有说明外，环$R$均为交换幺环。
%\begin{definition}
%真理想$P\subsetneq R$若满足
%\[
%ab\in P\Longrightarrow a\in P\ \text{或}\ b\in P,
%\]
%则称为\textbf{素理想}。
%真理想$M\subsetneq R$若不存在理想$A$满足$M\subsetneq A\subsetneq R$，则称为\textbf{极大理想}。
%\end{definition}
%
%\begin{example}
%对素数$p$，$(p)$是$\Z$的素理想，因为$p\mid ab$蕴含$p\mid a$或$p\mid b$。它也极大：若$(p)\subseteq(n)\subseteq\Z$，则$n\mid p$，故$(n)=(p)$或$(n)=\Z$。
%\end{example}
%\begin{proof}[补充证明]
%说明$\Z$中其余非零真理想既非素理想也非极大理想。任意这样的理想写成$(n)$，其中$n>1$为合数。取分解$n=ab$，$1<a,b<n$。则
%\[
%ab\in(n),\qquad a,b\notin(n),
%\]
%故$(n)$不素。又因$a\mid n$且$1<a<n$，
%\[
%(n)\subsetneq(a)\subsetneq\Z,
%\]
%故$(n)$不极大。零理想$(0)$是素理想，因为$\Z$为整环，但它不极大。
%\end{proof}
%
%\begin{theorem}\label{p2:prime-maximal}
%设$R$为交换幺环。
%\begin{enumerate}
%\item 若$\sigma:R\to R'$为满环同态，$N=\ker\sigma$，则$R$中包含$N$的极大理想与$R'$的极大理想一一对应。
%\item $R$为域，当且仅当$(0)$为极大理想。
%\item 真理想$M$极大，当且仅当$R/M$为域。
%\end{enumerate}
%\end{theorem}
%\begin{proof}
%第一项由环的对应定理给出：该对应保持包含关系，因此也保持“真理想中极大”这一性质。
%
%若$R$为域，任何非零理想$H$包含某个$a\ne0$，于是$1=a^{-1}a\in H$，故$H=R$，$(0)$极大。
%反之，若$(0)$极大而$a\ne0$，则包含$a$的理想$Ra$非零，必等于$R$。因此$1=ba$对某个$b\in R$成立，$a$可逆，故$R$为域。
%
%最后将第一项应用于自然满同态$R\to R/M$，再用第二项即可。
%\end{proof}
%
%\begin{theorem}\label{p2:prime-domain}
%设$R$为交换幺环。
%\begin{enumerate}
%\item 真理想$P$为素理想，当且仅当$R/P$为整环。
%\item $R$为整环，当且仅当$(0)$为素理想。
%\item 满环同态$\sigma:R\to R'$诱导包含$\ker\sigma$的素理想与$R'$的素理想之间的一一对应。
%\end{enumerate}
%因此每个极大理想都是素理想。
%\end{theorem}
%\begin{proof}
%在$R/P$中，
%\[
%(a+P)(b+P)=0+P\iff ab\in P.
%\]
%故$P$的素性恰好说商环中两个元素乘积为零时至少一个为零。取$P=(0)$即得第二项。
%
%由环同态定理，若$\sigma(H)=H'$且$\ker\sigma\subseteq H$，则$R/H\cong R'/H'$。根据第一项，两商环同时为整环或同时不是整环，这就证明第三项。
%
%域是整环；由前一定理，极大理想的商环为域，故它是素理想。
%\end{proof}
%
%\begin{theorem}\label{p2:prime-avoid-powers}
%若$R$为交换幺环，$a\in R$为非幂零元素，则存在素理想$P$，使得$a^m\notin P$对所有$m\geq0$成立。
%\end{theorem}
%\begin{proof}
%令$S$为所有与$\{a^m:m\geq0\}$不相交的理想组成的集合。因$a$非幂零，$(0)\in S$。
%
%按包含关系排序。若$\{A_\alpha\}$为$S$中的链，则其并$A=\bigcup_\alpha A_\alpha$为理想：任取$b,c\in A$，两者所在的两个理想可比较，故$b-c$属于较大的那个；$rb\in A$同理。并集仍避开所有$a^m$。由佐恩引理，$S$有极大元$P$。
%
%证明$P$为素理想。若$bc\in P$但$b,c\notin P$，则理想$H=Rb+P$、$K=Rc+P$都真包含$P$，因而都不属于$S$。于是有$a^r\in H$、$a^s\in K$。但
%\[
%HK=(Rb+P)(Rc+P)\subseteq P,
%\]
%因为$bc\in P$且其余各项均含$P$。故$a^{r+s}\in P$，矛盾。于是$P$为素理想。
%\end{proof}
%
%\begin{corollary}
%每个非零交换幺环至少有一个极大理想。令$r(R)$为全部素理想之交，则$r(R)$恰为$R$中全体幂零元素的集合，称为$R$的\textbf{诣零根}。
%\end{corollary}
%\begin{proof}
%在定理\ref{p2:prime-avoid-powers}的证明中取$a=1$。因为$R$非零，$1$不是幂零元素，而不含$1$的理想恰好是真理想。证明中所取的$P$在全体真理想中极大，故是$R$的极大理想。
%
%若$a^m=0$，则在任意素理想$P$中有$a^m\in P$。反复使用素性得$a\in P$，故$a\in r(R)$。
%
%若$a$不是幂零元素，上一定理给出一个不含$a$的素理想，因此$a\notin r(R)$。两边结合即得结论。
%\end{proof}
%
%\begin{example}
%$\Z$的诣零根为$(0)$。若$n=\prod_i p_i^{e_i}$，则$\Z/(n)$的素理想对应于$\Z$中包含$(n)$的素理想，即各$(p_i)/(n)$。因此
%\[
%r(\Z/(n))=\bigcap_i(p_i)/(n)=\bigl(\prod_i p_i\bigr)/(n).
%\]
%\end{example}
%
%\section{商域和分式环}
%
%\begin{definition}
%设$R$为整环。域$F$称为$R$的\textbf{商域}，若$R$可视为$F$的子环，且$F$的每个元素均可写成$a/b$，其中$a,b\in R$、$b\ne0$。
%\end{definition}
%
%\begin{theorem}\label{p2:fraction-field}
%每个整环都有商域。
%\end{theorem}
%\begin{proof}
%令$R^*=R\setminus\{0\}$，在$T=R\times R^*$上定义
%\[
%(a,b)\sim(c,d)\iff ad=bc.
%\]
%自反、对称显然。若$(a,b)\sim(c,d)$且$(c,d)\sim(e,f)$，则$ad=bc$、$cf=de$，从而
%\[
%adf=bcf=bde.
%\]
%因$d\ne0$且$R$为整环，消去$d$得$af=be$，故关系传递。
%
%将$(a,b)$的等价类记为$a/b$，定义
%\[
%\frac ab+\frac cd=\frac{ad+bc}{bd},\qquad
%\frac ab\frac cd=\frac{ac}{bd}.
%\]
%若$a/b=a'/b'$、$c/d=c'/d'$，即$ab'=a'b$、$cd'=c'd$，交叉相乘可得
%\begin{align*}
%(ad+bc)b'd'&=(a'd'+b'c')bd,\\
%(ac)b'd'&=(a'c')bd.
%\end{align*}
%故运算与代表元选择无关。分配律等环公理由$R$的相应等式逐项得到。$0/b$与$b/b$分别为零元和单位元；当$a\ne0$时，$b/a$是$a/b$的逆元。因此等价类集合$F$为域。
%
%映射$R\to F$，$a\mapsto a/1$为单环同态，且每个元素$a/b=(a/1)(b/1)^{-1}$，故$F$是商域。
%
%\end{proof}
%
%\begin{theorem}\label{p2:fraction-field-extension}
%设$F$为整环$R$的商域，$F'$为一个域。任一单环同态$\eta:R\to F'$都唯一延拓为单环同态$\widehat\eta:F\to F'$。
%\end{theorem}
%\begin{proof}
%先定义
%\[
%\widehat\eta:F\to F',\qquad \frac ab\longmapsto\eta(a)\eta(b)^{-1}.
%\]
%若$a/b=c/d$，则$ad=bc$，施加$\eta$得到$\eta(a)\eta(d)=\eta(b)\eta(c)$。由于$\eta(b),\eta(d)\ne0$，公式与分式的代表元选择无关。
%
%分式的加法和乘法公式表明$\widehat\eta$保运算，且$\widehat\eta(a/1)=\eta(a)$。若$\widehat\eta(a/b)=0$，则$\eta(a)=0$，由$\eta$单射得$a=0$，故$\widehat\eta$也单射。
%
%若$\eta':F\to F'$也是$\eta$的延拓，由$b(b^{-1})=1$及同态性可知$\eta'(b^{-1})=\eta(b)^{-1}$，于是$\eta'(a/b)=\eta(a)\eta(b)^{-1}$。故延拓唯一。
%\end{proof}
%
%\begin{corollary}
%整环的商域在保持原环元素不变的同构意义下唯一。
%\end{corollary}
%\begin{proof}
%若$F'$也是$R$的商域，将上一定理应用于包含映射$R\hookrightarrow F'$。所得$\widehat\eta:F\to F'$为单同态；$F'$中每个元素都是两个$R$元素的商，故又是满射。于是它为所需同构。
%\end{proof}
%
%\begin{definition}
%设$R$为交换幺环，$S\subseteq R$为非空的\textbf{乘性子集}，即$s,t\in S$蕴含$st\in S$。记
%\[
%N=\{a\in R:\text{存在 }s\in S\text{ 使 }as=0\}.
%\]
%若$0\notin S$，则$N$为$R$的理想，且$N\cap S=\varnothing$。
%\end{definition}
%\begin{proof}[补充证明]
%若$a,b\in N$，取$s,t\in S$使$as=bt=0$。于是
%\[
%(a-b)st=0,
%\]
%故$a-b\in N$；对$r\in R$，$(ra)s=0$，故$ra\in N$。若$a\in N\cap S$，又有$as=0$且$a,s\in S$，则$0\in S$，矛盾。
%\end{proof}
%
%\begin{definition}
%假设$0\notin S$。在$T=R\times S$上规定
%\[
%(a,s)\sim(b,t)\iff at-bs\in N.
%\]
%等价类记为$a/s$。集合$S^{-1}R=T/\!\sim$按
%\[
%\frac as+\frac bt=\frac{at+bs}{st},\qquad
%\frac as\frac bt=\frac{ab}{st}
%\]
%构成环，称为$R$关于$S$的\textbf{分式环}。
%\end{definition}
%\begin{proof}[补充证明]
%自反、对称由$N$的加法封闭性得到。设$(a,s)\sim(b,t)$、$(b,t)\sim(c,u)$。因为$N$为理想，
%\[
%u(at-bs)+s(bu-ct)=t(au-cs)\in N.
%\]
%由$N$的定义，存在$v\in S$使$vt(au-cs)=0$；$vt\in S$，故$au-cs\in N$，关系传递。
%
%加法与乘法的代表元无关性可由同样的交叉相乘得到。例如若$a/s=a'/s'$且$b/t=b'/t'$，则
%\begin{align*}
%(at+bs)s't'-(a't'+b's')st
%&=tt'(as'-a's)+ss'(bt'-b't)\in N.
%\end{align*}
%乘法的相应差也由$as'-a's$及$bt'-b't$的倍数组成，故属于$N$。环公理由$R$的环公理逐项得到。任取$s\in S$，$0/s$是零元、$s/s$是单位元，且不同$s$给出同一个单位元。
%\end{proof}
%
%\begin{theorem}\label{p2:localization-existence}
%设$0\notin S$。有自然环同态$\sigma:R\to S^{-1}R$，其核为$N$，每个$\sigma(s)$（$s\in S$）可逆，且$S^{-1}R$中的每个元素都可写成$\sigma(a)\sigma(s)^{-1}$。
%\end{theorem}
%\begin{proof}
%任取一个固定的$t\in S$，定义$\sigma(a)=at/t$。更换$t$所得的类相同，因为
%\[
%(at)u-(au)t=0.
%\]
%加法、乘法公式说明$\sigma$为环同态。又
%\[
%\sigma(a)=0\iff at/t=0/t\iff at^2\in N\iff a\in N.
%\]
%最后一个等价式利用$N$的定义与$S$的乘法封闭性。对$s\in S$，$\sigma(s)$的逆元为$t/(st)$，由乘法公式可核对。直接计算给出
%\[
%\frac as=\sigma(a)\sigma(s)^{-1}.
%\]
%\end{proof}
%
%\begin{theorem}\label{p2:localization-universal}
%若$\eta:R\to R'$为交换幺环之间的同态，$\eta(1)=1'$，且每个$\eta(s)$（$s\in S$）在$R'$中可逆，则存在唯一同态$\xi:S^{-1}R\to R'$使$\eta=\xi\circ\sigma$，并且
%\[
%\xi\!\left(\frac as\right)=\eta(a)\eta(s)^{-1}.
%\]
%\end{theorem}
%\begin{proof}[补充证明]
%若$a/s=b/t$，则$at-bs\in N$，取$u\in S$使$u(at-bs)=0$。施加$\eta$后，$\eta(u)$可逆，故
%\[
%\eta(a)\eta(t)=\eta(b)\eta(s),
%\]
%所列公式良定义。分式运算公式给出它保加法和乘法，且$\xi(\sigma(a))=\eta(a)$。
%
%任意满足$\eta=\xi\circ\sigma$的同态必须满足
%\[
%\xi(a/s)=\xi(\sigma(a))\xi(\sigma(s))^{-1}
%=\eta(a)\eta(s)^{-1},
%\]
%故唯一。
%\end{proof}
%
%\begin{example}
%对素数$p$，取$R=\Z$、$S=\Z\setminus(p)$。则$S^{-1}\Z$由分母不被$p$整除的有理数组成，它的非单位元恰为$p$的倍数；唯一的非零素理想是$pS^{-1}\Z$。
%\end{example}
%
%\section{交换环上的多项式环}
%
%设$R$为交换幺环。若$R$是扩环$R'$的子环，且两环有相同的单位元素，则可在$R'$中讨论由$R$与$u\in R'$生成的子环$R[u]$。它由有限和
%\[
%a_0+a_1u+\cdots+a_nu^n,\qquad a_i\in R,
%\]
%组成。为了不依赖$u$可能满足的代数关系，下面先构造形式多项式。
%
%\begin{definition}
%用$R[[x]]$表示全部无限序列$(a_0,a_1,\ldots)$，$a_i\in R$。定义
%\[
%(a_i)+(b_i)=(a_i+b_i),\qquad
%(a_i)(b_i)=(c_i),\qquad c_n=\sum_{i+j=n}a_ib_j.
%\]
%它称为$R$上的\textbf{一元形式幂级数环}。
%\end{definition}
%\begin{proof}[补充证明]
%每个$c_n$只涉及$n+1$项。乘法结合律在第$n$个系数上化为
%\[
%\sum_{i+j+k=n}(a_ib_j)c_k
%=\sum_{i+j+k=n}a_i(b_jc_k).
%\]
%交换律与分配律同样按系数验证。零元为$(0,0,\ldots)$，单位元为$(1,0,\ldots)$，故它是交换幺环。
%\end{proof}
%
%\begin{definition}
%$R[[x]]$中仅有有限个非零系数的序列组成子环$R[x]$，称为$R$上的\textbf{一元多项式环}。把$(a,0,\ldots)$与$a\in R$等同，记$x=(0,1,0,\ldots)$。于是每个多项式唯一写成
%\[
%f(x)=a_0+a_1x+\cdots+a_nx^n.
%\]
%\end{definition}
%\begin{proof}[补充证明]
%两个最终为零的序列之和仍最终为零。若$a_i=0$对$i>m$成立，$b_j=0$对$j>n$成立，则乘积的第$k$个系数在$k>m+n$时为零，故$R[x]$对乘法封闭。
%
%按系数乘法可知$x^i$只在第$i$个位置有系数$1$，故$\sum_i a_ix^i$正是序列$(a_0,a_1,\ldots)$。两个序列相等当且仅当每个系数相等，所以表示唯一。
%\end{proof}
%
%\begin{theorem}\label{p2:poly-universal}
%设$\sigma:R\to S$为交换幺环同态且$\sigma(1)=1'$，$u\in S$。则存在唯一环同态$\sigma_u:R[x]\to S$，其在$R$上的限制为$\sigma$，且$\sigma_u(x)=u$。具体地，
%\[
%\sigma_u\!\left(\sum_{i=0}^n a_ix^i\right)
%=\sum_{i=0}^n\sigma(a_i)u^i.
%\]
%其像为$\sigma(R)[u]$。
%\end{theorem}
%\begin{proof}
%多项式的系数表示唯一，因此公式定义了映射。它显然保加法。若$f=\sum_i a_ix^i$、$g=\sum_jb_jx^j$，则
%\begin{align*}
%\sigma_u(fg)
%&=\sum_k\sigma\!\left(\sum_{i+j=k}a_ib_j\right)u^k\\
%&=\sum_{i,j}\sigma(a_i)\sigma(b_j)u^{i+j}\\
%&=\sigma_u(f)\sigma_u(g).
%\end{align*}
%又$\sigma_u(a)=\sigma(a)$、$\sigma_u(x)=u$。由于$R[x]$由$R$和$x$生成，任何满足这些条件的同态在每个$\sum_i a_ix^i$上的值都已确定，故唯一。
%\end{proof}
%
%\begin{corollary}
%若扩环$S\supseteq R$与$R$有相同的单位元素，$u\in S$，则存在理想$I\subseteq R[x]$使
%\[
%R[u]\cong R[x]/I,\qquad I\cap R=(0).
%\]
%\end{corollary}
%\begin{proof}
%将上一定理用于包含映射$R\hookrightarrow S$。所得代入同态的像为$R[u]$，核为$I$。由第一同构定理得所列商环同构；因该同态在$R$上单射，$I\cap R=(0)$。
%\end{proof}
%
%\begin{definition}
%递归地定义$R[x_1,\ldots,x_n]=R[x_1,\ldots,x_{n-1}][x_n]$，称为$R$上的\textbf{$n$元多项式环}。它的每个元素唯一写成有限和
%\[
%f(x_1,\ldots,x_n)=\sum_{i_1,\ldots,i_n\geq0}
%a_{i_1\cdots i_n}x_1^{i_1}\cdots x_n^{i_n},
%\qquad a_{i_1\cdots i_n}\in R.
%\]
%\end{definition}
%\begin{proof}[补充证明]
%对$n$归纳。一元情形已证。若$n-1$元时成立，把$f$看成关于$x_n$的多项式
%\[
%f=\sum_{j=0}^m f_j(x_1,\ldots,x_{n-1})x_n^j.
%\]
%各$f_j$由归纳假设唯一展开，代入即得有限多项式和。若$f=0$，则每个$f_j=0$，再由归纳假设可知全部$a_{i_1\cdots i_n}=0$。
%\end{proof}
%
%\begin{theorem}\label{p2:multi-poly-universal}
%设$\sigma:R\to S$为保持单位元素的环同态，$u_1,\ldots,u_n\in S$。则$\sigma$唯一延拓为环同态
%\[
%\sigma_{\mathbf u}:R[x_1,\ldots,x_n]\to S,
%\qquad x_i\longmapsto u_i,
%\]
%其像为$\sigma(R)[u_1,\ldots,u_n]$。
%\end{theorem}
%\begin{proof}
%依次对$x_1,\ldots,x_n$应用一元多项式的延拓定理。每一步的系数环是前面变量构成的多项式环，所取$u_i$与先前各像交换，因为$S$交换。所得映射将有限和逐项代入，故其像是所述生成子环。多项式由$R,x_1,\ldots,x_n$生成，因而延拓唯一。
%\end{proof}
%
%\begin{corollary}
%若$S\supseteq R$为同单位的扩环，$u_1,\ldots,u_n\in S$，则对某理想$I\subseteq R[x_1,\ldots,x_n]$有
%\[
%R[u_1,\ldots,u_n]\cong R[x_1,\ldots,x_n]/I,
%\qquad I\cap R=(0).
%\]
%\end{corollary}
%\begin{proof}
%取代入同态的核$I$，应用第一同构定理。由于代入映射在$R$上是包含映射，其核与$R$相交为零。
%\end{proof}
%
%\begin{corollary}
%若$N$为$R$的理想，则$N[x_1,\ldots,x_n]$为$R[x_1,\ldots,x_n]$的理想，且
%\[
%R[x_1,\ldots,x_n]/N[x_1,\ldots,x_n]
%\cong (R/N)[x_1,\ldots,x_n].
%\]
%\end{corollary}
%\begin{proof}
%自然满同态$R\to R/N$逐系数延拓，得到满同态
%\[
%\sum a_{\mathbf i}x^{\mathbf i}
%\longmapsto\sum(a_{\mathbf i}+N)x^{\mathbf i}.
%\]
%目标多项式为零，当且仅当各$a_{\mathbf i}\in N$，故核正是$N[x_1,\ldots,x_n]$。应用第一同构定理即得结论。
%\end{proof}
%
%\section{整环上的一元多项式环}
%
%设$R$为整环。非零多项式$f(x)=a_nx^n+\cdots+a_0$（$a_n\ne0$）的\textbf{次数}为$n$，记作$\deg f=n$；约定$\deg0=-\infty$。
%
%\begin{proposition}
%对$f,g\in R[x]$，
%\[
%\deg(f+g)\leq\max\{\deg f,\deg g\},\qquad
%\deg(fg)=\deg f+\deg g\quad(f,g\ne0).
%\]
%因而$R[x]$也是整环，且它的单位恰为$R$的单位。
%\end{proposition}
%\begin{proof}[补充证明]
%加法的最高非零项不可能超过原来两多项式的最高次数。设$f,g$的首项系数分别为$a_m,b_n$。乘积的$x^{m+n}$系数为$a_mb_n\ne0$，因为$R$无零因子，故次数相加。
%
%若$fg=1$，则$\deg f+\deg g=0$，故$f,g$都是非零常数，且在$R$中互为逆元。反向显然。
%\end{proof}
%
%\begin{theorem}\label{p2:poly-division}
%\textbf{除法算式：}设$R$为交换幺环，$f,g\in R[x]$，$g\ne0$且$g$的首项系数可逆。则唯一存在$q,r\in R[x]$使
%\[
%f=qg+r,\qquad r=0\ \text{或}\ \deg r<\deg g.
%\]
%\end{theorem}
%\begin{proof}[补充证明]
%对$\deg f$归纳。若$f=0$或$\deg f<\deg g$，取$q=0,r=f$。否则设$f$首项为$ax^m$、$g$首项为$bx^n$，其中$b$为单位，且$m\geq n$。令
%\[
%f_1=f-ab^{-1}x^{m-n}g.
%\]
%首项相消，故$\deg f_1<m$。由归纳假设将$f_1$除以$g$，得到所需$q,r$。
%
%若$f=qg+r=q'g+r'$，则$(q-q')g=r'-r$。若$q\ne q'$，其左边次数至少为$\deg g$，而右边次数小于$\deg g$，矛盾。故$q=q'$、$r=r'$。
%\end{proof}
%
%\begin{corollary}
%对$c\in R$、$f\in R[x]$，有
%\[
%f(x)=q(x)(x-c)+f(c).
%\]
%特别地，$(x-c)\mid f(x)$当且仅当$f(c)=0$。
%\end{corollary}
%\begin{proof}[补充证明]
%用除法算式除以首项系数为$1$的$x-c$，余式是常数$r$。代入$c$得$r=f(c)$，故结论成立。
%\end{proof}
%
%\begin{theorem}\label{p2:poly-roots}
%设$R$为整环，$f\in R[x]$为$n$次非零多项式，则$f$在$R$中至多有$n$个不同的根。
%\end{theorem}
%\begin{proof}
%把$R$嵌入其商域$F$，在$F[x]$中看待$f$。对$n$归纳。若$c$为根，则$f=(x-c)g$，$\deg g=n-1$。任意另一根$d\ne c$满足
%\[
%0=f(d)=(d-c)g(d).
%\]
%在域$F$中$d-c\ne0$，故$g(d)=0$。由归纳假设，除$c$外至多还有$n-1$个根。
%\end{proof}
%
%\begin{remark}
%整环条件不能省去。例如在$\Z/(8)$中，$x^2-1$有$1,3,5,7$四个根。
%\end{remark}
%
%\begin{theorem}\label{p2:finite-multiplicative-subgroup}
%若$R$为整环，则$R\setminus\{0\}$的任一有限乘法子群都是循环群。
%\end{theorem}
%\begin{proof}
%设$G$为这样的有限子群。由于$R$交换，$G$为有限交换群。令$m$为$G$中元素阶的最小公倍数。有限交换群的循环群结论给出一个元素的阶为$m$，并且每个元素满足$x^m=1$。
%
%多项式$x^m-1$在整环$R$中至多有$m$个根，故$|G|\leq m$；但阶为$m$的元素生成$m$个不同元素，故$|G|=m$，$G$循环。
%\end{proof}
%
%\begin{corollary}
%若$F$为含$q$个元素的有限域，则$F^\times$为$q-1$阶循环群，且$F$中的每个元素均满足$x^q=x$。
%\end{corollary}
%\begin{proof}
%$F^\times$有$q-1$个元素，故由上一定理循环。对非零$x$，拉格朗日定理给出$x^{q-1}=1$，所以$x^q=x$；零元也满足等式。
%\end{proof}
%
%\begin{definition}
%交换幺环$R$中，每个理想均为主理想的整环称为\textbf{主理想整环}。
%\end{definition}
%\begin{theorem}\label{p2:field-poly-pid}
%若$F$为域，则$F[x]$为主理想整环。
%\end{theorem}
%\begin{proof}
%设$I\ne(0)$为$F[x]$的理想，取其中次数最小的非零多项式$f$。因$F$为域，可将$f$乘非零常数，使首项系数为$1$。
%
%显然$(f)\subseteq I$。任取$g\in I$，用除法算式写成
%\[
%g=qf+r,\qquad r=0\ \text{或}\ \deg r<\deg f.
%\]
%因$I$为理想，$r=g-qf\in I$。最小次数的选择迫使$r=0$，故$g\in(f)$，于是$I=(f)$。零理想由$0$生成。
%\end{proof}
%
%\begin{definition}
%在整环$R$中，非零且非单位的元素$f$若不能写成两个非单位的乘积，则称为\textbf{不可约元素}。在$F[x]$中，它们就是不可约多项式。
%\end{definition}
%\begin{theorem}\label{p2:irreducible-maximal}
%设$F$为域，$f\in F[x]$且$\deg f\geq1$。以下条件等价：
%\begin{enumerate}
%\item $f$不可约；
%\item $(f)$为极大理想；
%\item $F[x]/(f)$为域；
%\item $F[x]/(f)$为整环；
%\item $(f)$为素理想。
%\end{enumerate}
%\end{theorem}
%\begin{proof}
%若$f$不可约且$(f)\subseteq I\subseteq F[x]$，主理想性质给出$I=(g)$，所以$g\mid f$。不可约性迫使$g$为单位或与$f$相伴，因此$I=F[x]$或$I=(f)$，故$(f)$极大。
%
%极大理想的商环为域；域为整环；商环为整环当且仅当$(f)$为素理想。这证明前四项及第五项的顺向关系。
%
%若$(f)$为素理想而$f=gh$，其中$g,h$均为非单位，则$gh\in(f)$，素性迫使$g\in(f)$或$h\in(f)$。例如$g\in(f)$意味着$f\mid g$，与$\deg g<\deg f$矛盾。故$f$不可约。
%\end{proof}
%
%\begin{corollary}
%若$S\supseteq F$为同单位的扩环，$u\in S$为$F$上的代数元，且代入同态$F[x]\to F[u]$的核由首项系数为$1$的多项式$f$生成，则$F[u]$为域当且仅当$f$不可约。
%\end{corollary}
%\begin{proof}
%由第一同构定理，$F[u]\cong F[x]/(f)$。结论随前一定理直接得到。
%\end{proof}
%
%\begin{example}
%若有限域$F$的特征为素数$p$，则素域$\F_p\subseteq F$。取$F^\times$的生成元$\alpha$。由于$F$的非零元素都是$\alpha$的幂，
%\[
%F=\F_p[\alpha].
%\]
%考虑代入同态$\sigma:\F_p[x]\to F$，$f(x)\mapsto f(\alpha)$。它满射，核$I$是$\F_p[x]$的理想。由定理\ref{p2:field-poly-pid}，$I=(f)$，其中$f$可取为首项系数为$1$的多项式。因$F$为域，由定理\ref{p2:irreducible-maximal}，$f$不可约。若$\deg f=n$，则
%\[
%F\cong\F_p[x]/(f),\qquad |F|=p^n.
%\]
%后一个等式来自商环中每个元素均有唯一的次数小于$n$的代表元，故有$p^n$个元素。
%\end{example}
%
%\begin{example}
%设$F$为域，$A\in M_n(F)$，把标量$a\in F$与纯量矩阵$aE$等同。则$F[A]$为由$F$与$A$生成的交换子环。由于$M_n(F)$作为$F$上向量空间的维数为$n^2$，$E,A,\ldots,A^{n^2}$线性相关，故存在非零$f\in F[x]$使$f(A)=0$。
%
%当$n=2$、$F=\Q$且$A$不是纯量矩阵时，$A$的极小多项式次数为$2$，记为$f$。若$f$不可约，则$\Q[A]\cong\Q[x]/(f)$为二次域。
%
%若$f=(x-x_1)(x-x_2)$，其中$x_1\ne x_2$，令
%\[
%B=\frac{A-x_1E}{x_2-x_1},\qquad
%C=\frac{A-x_2E}{x_1-x_2}.
%\]
%由$f(A)=0$得$BC=0$，并有$B+C=E$、$B^2=B$、$C^2=C$。因此
%\[
%\Q[A]=\Q B\oplus\Q C.
%\]
%若$f=(x-x_1)^2$，令$B=A-x_1E$，则$B\ne0$且$B^2=0$，故$\Q[A]$含非零零幂元。
%\end{example}
%
%\section{多项式函数}
%
%设$F$为域。函数集合$F^F$按逐点加法、乘法构成交换幺环。每个$a\in F$对应常数函数$f_a(s)=a$，由此把$F$视为$F^F$的子环。恒等函数$s\mapsto s$记为自变量函数$s$；由它与常数函数生成的子环$F[s]$称为$F$上的\textbf{多项式函数环}。
%\begin{proof}[补充验证]
%对任意函数$f,g,h:F\to F$及$s\in F$，有
%\[
%((fg)h)(s)=(f(s)g(s))h(s)=f(s)(g(s)h(s))=(f(gh))(s).
%\]
%加法结合律、交换律及分配律也逐点归结为$F$中的相应等式。零函数和常数函数$1$分别为零元和单位元，因此$F^F$为交换幺环。
%\end{proof}
%
%\begin{theorem}\label{p2:poly-functions-one}
%设$F$为域，$s:F\to F$为恒等函数。
%\begin{enumerate}
%\item 若$F$为无限域，则代入映射$F[x]\to F[s]$为同构；多项式与多项式函数一一对应。
%\item 若$F$有$q$个元素，则
%\[
%F[s]\cong F[x]/(x^q-x)=F^F.
%\]
%每个多项式函数由唯一的次数小于$q$的多项式表示。
%\end{enumerate}
%\end{theorem}
%\begin{proof}
%由多项式环的延拓定理，代入映射$\operatorname{ev}:F[x]\to F^F$的像为$F[s]$。
%
%若$F$无限而$f\in\ker\operatorname{ev}$，则$f(a)=0$对所有$a\in F$成立。非零多项式只有有限多个根，所以$f=0$，故代入映射单射。
%
%若$|F|=q$，先取所有元素对应的一次因子的乘积
%\[
%g(x)=\prod_{a\in F}(x-a).
%\]
%它是$q$次首项系数为$1$的多项式。若$f\in\ker\operatorname{ev}$，则$F$中每个$a$都是$f$的根。由因式定理及根互不相同，$g\mid f$；因此$\ker\operatorname{ev}=(g)$。
%
%另一方面，有限域的每个元素满足$a^q=a$，所以$g\mid(x^q-x)$。两者都是$q$次首项系数为$1$的多项式，故$g=x^q-x$，从而
%\[
%\ker\operatorname{ev}=(x^q-x).
%\]
%由除法算式，商环中的元素由次数小于$q$的多项式唯一表示，共有$q^q$个。函数集合$F^F$也有$q^q$个元素，因此代入映射满到$F^F$。
%\end{proof}
%
%\begin{definition}
%对$n\geq1$，令$F^{(n)}=F^n$，在函数环$F^{F^{(n)}}$中定义坐标函数
%\[
%s_i(a_1,\ldots,a_n)=a_i\qquad(1\leq i\leq n).
%\]
%由$F$的常数函数与$s_1,\ldots,s_n$生成的子环$F[s_1,\ldots,s_n]$称为$F$上的\textbf{$n$元多项式函数环}。
%\end{definition}
%
%\begin{theorem}\label{p2:poly-functions-many}
%设$F$为域，代入映射$\operatorname{ev}:F[x_1,\ldots,x_n]\to F[s_1,\ldots,s_n]$由$x_i\mapsto s_i$给出。
%\begin{enumerate}
%\item 若$F$无限，则$\operatorname{ev}$为同构。
%\item 若$|F|=q$，则
%\[
%\ker\operatorname{ev}=(x_1^q-x_1,\ldots,x_n^q-x_n),
%\qquad F[s_1,\ldots,s_n]=F^{F^{(n)}}.
%\]
%每个函数由唯一一个各变量次数均小于$q$的多项式表示。
%\end{enumerate}
%\end{theorem}
%\begin{proof}
%代入映射的像按定义为$F[s_1,\ldots,s_n]$。
%
%若$F$无限，对$n$归纳。设$f$在$F^n$上恒为零，把它写成
%\[
%f(x_1,\ldots,x_n)
%=\sum_{j=0}^m f_j(x_1,\ldots,x_{n-1})x_n^j.
%\]
%固定前$n-1$个变量的任意取值，右边作为$x_n$的一元多项式在无限多点为零，故各系数值均为零。由归纳假设，每个$f_j$都是零多项式，因而$f=0$。
%
%若$F$有$q$个元素，记$I=(x_1^q-x_1,\ldots,x_n^q-x_n)$。每个生成元在$F^n$上恒为零，所以$I\subseteq\ker\operatorname{ev}$。
%
%先对任意$f\in F[x_1,\ldots,x_n]$，依次以首项系数为$1$的$x_i^q-x_i$作关于$x_i$的除法。由此可写
%\[
%f=\sum_{i=1}^n q_i(x_1,\ldots,x_n)(x_i^q-x_i)
%+r(x_1,\ldots,x_n),
%\]
%其中$r$对每个$x_i$的次数都小于$q$。若原多项式在$F^n$上恒为零，则$r$也恒为零。
%
%对$n$归纳证明这样的$r$只能是零多项式。当$n=1$时，这是根的个数定理。设$n>1$，把$r$按$x_n$展开为
%\[
%r=\sum_{j=0}^{q-1}r_j(x_1,\ldots,x_{n-1})x_n^j.
%\]
%固定前$n-1$个变量的任意取值，再让$x_n$遍历$F$。右边是次数小于$q$、却有$q$个根的一元多项式，故每个$r_j$在$F^{n-1}$上恒为零。归纳假设给出$r_j=0$，从而$r=0$。因此$\ker\operatorname{ev}=I$。
%
%各变量次数均小于$q$的多项式共有$q^{q^n}$个，函数集合$F^{F^n}$也有$q^{q^n}$个元素。代入映射在商环上单射，且两集合有限等势，所以满射。
%\end{proof}
%\section{章末习题中的知识与解法}
%% !TEX root = ../../main.tex
% 第三章习题；由 chapter3.tex 按需输入。

\subsection*{环、理想与四元数（习题1--20）}

\begin{problem}
\textbf{习题1：}设环$R$中$1-ab$可逆，证明$1-ba$可逆。
\end{problem}
\begin{proof}[补充证明]
令$c=(1-ab)^{-1}$。由$c(1-ab)=(1-ab)c=1$，直接计算得到
\begin{align*}
(1-ba)(1+bca)&=1+b(c-1-abc)a=1,\\
(1+bca)(1-ba)&=1+b(c-1-cab)a=1.
\end{align*}
因此$(1-ba)^{-1}=1+b(1-ab)^{-1}a$。
\end{proof}

\begin{problem}
\textbf{习题2：}若$u$有右逆，证明以下条件等价：$u$有多于一个右逆；$u$是零因子；$u$不是单位。
\end{problem}
\begin{proof}[补充证明]
取$uv=1$。若$uw=1$且$w\ne v$，则$u(w-v)=0$，所以$u$是左零因子。若$u$是零因子，由$uv=1$可知$xu=0$蕴含$x=xuv=0$；于是存在非零$t$使$ut=0$，而$u(v+t)=1$，故右逆不唯一。

若$u$是单位，则右逆必为其唯一的逆。反之，若右逆$v$唯一，则
\[
u(vu-1)=0,
\]
故$u(v+vu-1)=1$。唯一性给出$vu=1$，从而$u$是单位。
\end{proof}

\begin{problem}
\textbf{习题3：}若$u$有多于一个右逆，证明它有无穷多个右逆。
\end{problem}
\begin{proof}[补充证明]
取$uv=1$。由于右逆不唯一，习题2说明$u$不是单位，故$e=1-vu\ne0$。又有$ue=0$和$ev=0$。对每个$k\geqslant0$，
\[
u(v+eu^k)=1.
\]
若$i<j$且$eu^i=eu^j$，右乘$v^j$，则左边为$ev^{j-i}=0$，右边为$e$，矛盾。因此这些右逆两两不同。
\end{proof}

\begin{problem}
\textbf{习题4：}证明交换环中全部幂零元构成理想，并判断非交换环中的情形。
\end{problem}
\begin{proof}[补充证明]
若$a^m=b^n=0$且$ab=ba$，则二项式展开中$(a+b)^{m+n-1}$的每一项均为零。任取$r\in R$，有$(ra)^m=r^ma^m=0$。故幂零元集对加法、取负和乘以环元封闭。

非交换情形不成立。例如在$M_2(\mathbb Z)$中，$E_{12}$和$E_{21}$均幂零，而$E_{12}+E_{21}$的平方为$1$。
\end{proof}

\begin{problem}
\textbf{习题5：}设$F$的特征为素数$p$，证明$(a+b)^p=a^p+b^p$。并将特征的讨论推广到幺环及不必有幺元的环。
\end{problem}
\begin{proof}[补充证明]
由于$F$交换，二项式公式和$p\mid\binom pk$（$0<k<p$）给出
\[
(a+b)^p=a^p+\sum_{k=1}^{p-1}\binom pk a^kb^{p-k}+b^p=a^p+b^p.
\]
幺环的特征是其加法群中$1$的阶，因而为$0$或某个正整数；若幺环是整环，正特征必须为素数。不带幺元的环以加法群的指数定义特征，即最小的使$nR=0$的正整数；若不存在则特征为$0$。交换且特征为$p$时，上述恒等式同样成立。
\end{proof}

\begin{problem}
\textbf{习题6：}设$I\lhd R$。证明$M_n(I)\lhd M_n(R)$，且$M_n(R)$的每个理想均如此得到。
\end{problem}
\begin{proof}[补充证明]
矩阵乘法按行列展开，$I$为双边理想立即给出$M_n(I)$为理想。

反过来，设$J\lhd M_n(R)$，令$I$为$J$中矩阵的全部分量所组成的集合。用矩阵单位可将$A=(a_{ij})\in J$的任一分量取出：
\[
E_{ki}AE_{j\ell}=a_{ij}E_{k\ell}\in J.
\]
于是$I$对加法及左右乘以$R$中元素封闭，且$aE_{ij}\in J$对所有$a\in I$成立。因此$J=M_n(I)$。对角元取出也说明$I$唯一。
\end{proof}

\begin{problem}
\textbf{习题7：}若$I,H,N$为$R$的理想，证明$(H+N)I=HI+NI$及$I(H+N)=IH+IN$。
\end{problem}
\begin{proof}[补充证明]
理想的积由有限个乘积之和构成。对每个$h\in H$、$n\in N$、$i\in I$，分配律给出$(h+n)i=hi+ni$，因此$(H+N)I\subseteq HI+NI$。反向包含由$H,N\subseteq H+N$得到。右侧等式同理。
\end{proof}

\begin{problem}
\textbf{习题8：}设$\eta:R\twoheadrightarrow R'$为保幺满同态。判断并证明或举反例：(i) 幂零元、幂等元的像仍是相应元素；(ii) 零因子的像仍是零因子；(iii) $R$为整环则$R'$为整环；(iv) $R'$为整环则$R$为整环；(v) 单位的像仍是单位；(vi) 像为单位的元素必为单位。
\end{problem}
\begin{proof}[补充证明]
(i) 对。若$a^m=0$，则$\eta(a)^m=0$；若$e^2=e$，则$\eta(e)^2=\eta(e)$。

(ii) 错。$\mathbb Z/6\mathbb Z\twoheadrightarrow\mathbb Z/3\mathbb Z$将零因子$4$送到非零单位$1$。

(iii) 错。$\mathbb Z\twoheadrightarrow\mathbb Z/6\mathbb Z$的源是整环，目标不是。

(iv) 错。$\mathbb Z/6\mathbb Z\twoheadrightarrow\mathbb Z/2\mathbb Z$的目标是域，源不是整环。

(v) 对。若$uv=vu=1$，则$\eta(u)\eta(v)=\eta(v)\eta(u)=1'$。

(vi) 错。$\mathbb Z\twoheadrightarrow\mathbb Z/5\mathbb Z$将非单位$2$送到单位。
\end{proof}

\begin{problem}
\textbf{习题9：}证明有理系数四元数$H_0=\{a_0+a_1I+a_2J+a_3K:a_i\in\mathbb Q\}$是$H$的子体。
\end{problem}
\begin{proof}[补充证明]
四元数的乘法表表明$H_0$对加法、乘法及取负封闭，且含$1$。若$0\ne q\in H_0$，则
\[
q^{-1}=\frac{\bar q}{q\bar q}
=\frac{a_0-a_1I-a_2J-a_3K}{a_0^2+a_1^2+a_2^2+a_3^2}\in H_0.
\]
分母是非零有理数，故$H_0$是子体。
\end{proof}

\begin{problem}
\textbf{习题10：}确定四元数体$H$的中心。
\end{problem}
\begin{proof}[补充证明]
设$q=a+bI+cJ+dK$与$I,J$都交换。比较$qI=Iq$的$J,K$系数，得$c=d=0$；再比较$qJ=Jq$，得$b=0$。故$q=a\in\mathbb R$。实数与所有四元数交换，因此中心为$\mathbb R$。
\end{proof}

\begin{problem}
\textbf{习题11：}求$H$中与$I$不交换的全部元素。
\end{problem}
\begin{proof}[补充证明]
对$q=a+bI+cJ+dK$，由$I J=K$、$J I=-K$、$I K=-J$、$K I=J$，
\[
Iq-qI=2cK-2dJ.
\]
所以所求为$\{a+bI+cJ+dK:c,d\text{不同时为零}\}$。
\end{proof}

\begin{problem}
\textbf{习题12：}设$S$是$H$的子体，且对所有$0\ne d\in H$有$dSd^{-1}\subseteq S$。证明$S=H$或$S$包含于中心。
\end{problem}
\begin{proof}[补充证明]
设$S\ne H$，取$x\in H\setminus S$及任意$s\in S$。由于$x,x+1$都非零，记
\[
a=xsx^{-1}\in S,\qquad b=(x+1)s(x+1)^{-1}\in S.
\]
由$xs=ax$及$xs+s=bx+b$得到$(a-b)x=b-s$。若$a\ne b$，则$x=(a-b)^{-1}(b-s)\in S$，矛盾。因此$a=b$，继而$s=b=a$，即$x$与$s$交换。

所以$S$外的每个元素都与$S$中的每个元素交换。对$y\in S$，$x+y\notin S$，于是$x+y$也与$s$交换；相减可知$y$与$s$交换。故$S$中的元素与$H$的每个元素交换，即$S$包含于中心。
\end{proof}

\begin{problem}
\textbf{习题13：}设$I,H,N\lhd R$，若$HN\subseteq I$且$I+H=R$，证明$N\subseteq I$。
\end{problem}
\begin{proof}[补充证明]
取$1=i+h$，其中$i\in I$、$h\in H$。任取$n\in N$，则
\[
n=(i+h)n=in+hn\in I.
\]
故$N\subseteq I$。
\end{proof}

\begin{problem}
\textbf{习题14：}设$I,H,N$为交换环$R$的理想，$I\subseteq H,N$且$H+N=R$，证明$HN\supseteq I$。
\end{problem}
\begin{proof}[补充证明]
取$1=h+n$，其中$h\in H$、$n\in N$。对每个$i\in I$，由于$I\subseteq H\cap N$，有
\[
i=ih+in\in HN.
\]
故$HN\supseteq I$。这里的包含可以是相等，例如$R=F\times F$、$I=0$、$H=F\times0$、$N=0\times F$。
\end{proof}

\begin{problem}
\textbf{习题15：}在$M_n(F)$（$n>1$）中找一个子环$R$，使$R$除恒等同构外没有其他反自同构。
\end{problem}
\begin{proof}[补充解答]
取$F$的素子域$F_0$，令$R=F_0\,1_n$。任一保幺自同构必须固定$1_n$，故固定由$1_n$生成的素子域；因$R$交换，反自同构与自同构相同。因此$R$仅有恒等反自同构。
\end{proof}

\begin{problem}
\textbf{习题16：}求$R=\left\{\begin{pmatrix}a&0\\b&c\end{pmatrix}:a,b,c\in\mathbb Z\right\}$的全部理想。
\end{problem}
\begin{proof}[补充解答]
全部理想恰为
\[
I_{m,k,n}=\left\{\begin{pmatrix}a&0\\b&c\end{pmatrix}:a\in m\mathbb Z,\ b\in k\mathbb Z,\ c\in n\mathbb Z\right\},
\]
其中$m,k,n\geqslant0$，并满足$m\mathbb Z,n\mathbb Z\subseteq k\mathbb Z$（即$k\mid m,n$，将$0$的整除意义理解为集合包含）。

事实上，用$E_{11},E_{22}$从理想中的矩阵左右相乘，可分别取出三个分量，故理想必为上述形式。再用$E_{21}$左右相乘，得到两个对角分量的取值都须属于左下角分量的取值集。反之，这两个包含关系保证上述集合对任意矩阵的左右乘法封闭。
\end{proof}

\begin{problem}
\textbf{习题17：}确定$M_2(\mathbb Z/(n))$的单位群的阶。
\end{problem}
\begin{proof}[补充解答]
由中国剩余定理，把$n=\prod p_i^{e_i}$分解，则
\[
\operatorname{GL}_2(\mathbb Z/(n))\cong\prod_i\operatorname{GL}_2(\mathbb Z/(p_i^{e_i})).
\]
模$p$约化到$\operatorname{GL}_2(\mathbb F_p)$，每个矩阵有$p^{4(e-1)}$个提升，而矩阵可逆当且仅当其模$p$像可逆。第一列有$p^2-1$种非零选择，第二列有$p^2-p$种不在第一列张成直线中的选择。故所求阶为
\[
\prod_{p^e\parallel n}p^{4(e-1)}(p^2-1)(p^2-p).
\]
当$n=1$时，该群只有一个元素。
\end{proof}

\begin{problem}
\textbf{习题18：}设交换环$R$不含零因子且仅有有限个非零元，证明$R$是有限除环。
\end{problem}
\begin{proof}[补充证明]
取$0\ne a\in R$。乘以$a$是有限集合$R$上的单射，因而满射；存在$e\in R$使$ae=a$。任取$b\in R$，由满射性写$b=ac$，于是$eb=eac=ac=b$。因此$e$是$R$的单位元。又对任意$0\ne b$，乘以$b$同样满射，存在$c$使$bc=e$。故$R$是有限域，特别是有限除环。
\end{proof}

\begin{problem}
\textbf{习题19：}若$N\lhd R$而$I\lhd N$，问$I$是否必为$R$的理想？
\end{problem}
\begin{proof}[补充解答]
未必。令$R$为$F$上所有严格上三角$3\times3$矩阵，$N=F E_{12}+F E_{13}$。直接由矩阵单位的乘法可知$N\lhd R$，且$N^2=0$，故$I=F E_{12}\lhd N$；但$E_{12}E_{23}=E_{13}\notin I$，所以$I$不是$R$的理想。
\end{proof}

\begin{problem}
\textbf{习题20：}设$R$为交换环、$a$不可逆，证明有极大理想包含$a$。
\end{problem}
\begin{proof}[补充证明]
由$a$不可逆，$(a)\ne R$。考虑所有包含$(a)$的真理想，按包含关系排序。任一全序子集的并仍是真理想：若它含$1$，则某个成员已含$1$，矛盾。由佐恩引理，有极大元$M$。它是$R$的极大理想，并含$a$。
\end{proof}

\subsection*{根、同态与素理想（习题21--32）}

\begin{problem}
\textbf{习题21：}设$J(R)$为交换环$R$的全部极大理想之交。证明$a\in J(R)$当且仅当$1-ax$对所有$x\in R$均可逆。
\end{problem}
\begin{proof}[补充证明]
若$a\in J(R)$而$1-ax$不是单位，则它属于某个极大理想$M$。因$a\in M$，便有$1=(1-ax)+ax\in M$，矛盾。

反之，若$a\notin J(R)$，取极大理想$M$使$a\notin M$。在域$R/M$中，$a+M$非零，故存在$x\in R$使$ax+M=1+M$。于是$1-ax\in M$，不是单位，矛盾。
\end{proof}

\begin{problem}
\textbf{习题22：}若$a$满足习题21的条件，称其为强拟正则元。证明全部强拟正则元构成$J(R)$。
\end{problem}
\begin{proof}[补充证明]
由习题21，对每个$a\in R$，
\[
a\text{强拟正则}\quad\Longleftrightarrow\quad a\in\bigcap_{M\text{极大}}M=J(R).
\]
故两集合相同。
\end{proof}

\begin{problem}
\textbf{习题23：}证明有限交换环的每个素理想都是极大理想。
\end{problem}
\begin{proof}[补充证明]
若$P$为素理想，则$R/P$是有限整环。由习题18，它是域，所以$P$是极大理想。
\end{proof}

\begin{problem}
\textbf{习题24：}证明第二章第2节的定理4和定理5等价。
\end{problem}
\begin{proof}[补充证明]
定理4说$\mathbb Z$的每个子群都是$m\mathbb Z$；定理5说无限循环群与$\mathbb Z$同构，并由此描述全部循环群的子群。

由定理4，若$G=\langle g\rangle$，满同态$\varphi:\mathbb Z\to G$，$k\mapsto g^k$，把$G$的任一子群$H$的原像写成$m\mathbb Z$。于是$H=\langle g^m\rangle$；当$G$无限时$\ker\varphi=0$，故$G\cong\mathbb Z$。若$G$的阶为$n$，则$\ker\varphi=n\mathbb Z\subseteq m\mathbb Z$，所以$m\mid n$，且$|H|=n/m$。反过来，$n$的每个正因子$m$都给出唯一的子群$\langle g^m\rangle$。这给出定理5。

反过来，把定理5应用于无限循环群$(\mathbb Z,+)=\langle1\rangle$，其任一子群由某个整数$m$生成，即为$m\mathbb Z$。取$m\geqslant0$并用子群的包含关系，即得定理4。
\end{proof}

\begin{problem}
\textbf{习题25：}设$G$为复数域中全部单位根构成的乘法群。给出一个满而非单的自同态。
\end{problem}
\begin{proof}[补充解答]
取$\varphi(z)=z^2$。若$z^m=1$，则$\varphi(z)^m=1$，故$\varphi$是自同态。对任意$m$次单位根$w$，取它的一个平方根$z$，则$z^{2m}=1$，所以$z\in G$且$\varphi(z)=w$。但$\varphi(1)=\varphi(-1)=1$，它并非单射。
\end{proof}

\begin{problem}
\textbf{习题26：}设$\eta:R\to R'$满足加法同态、$\eta(1)=1'$及$\eta(aba)=\eta(a)\eta(b)\eta(a)$。证明$\eta$为若尔当同态，并推出幂与极化恒等式。
\end{problem}
\begin{proof}[补充证明]
在给定恒等式中取$b=1$，即得$\eta(a^2)=\eta(a)^2$。加法性与保幺性均已给出，因此$\eta$为若尔当同态。取三重乘积中的两个外侧变量都为$a$、中间变量为$a^m$，得到
\[
\eta(a^{m+2})=\eta(a)\eta(a^m)\eta(a).
\]
以$k=0,1$为起点分别对偶数和奇数归纳，得$\eta(a^k)=\eta(a)^k$。将$a$换成$a+b$并减去平方项，得
\[
\eta(ab+ba)=\eta(a)\eta(b)+\eta(b)\eta(a).
\]
对$abc+cba$，把$a+c$代入三重乘积恒等式并减去$a$、$c$各自的项，得到
\[
\eta(abc+cba)=\eta(a)\eta(b)\eta(c)+\eta(c)\eta(b)\eta(a).
\]
\end{proof}

\begin{problem}
\textbf{习题27：}设$\eta:R\to R'$是到无零因子环的若尔当同态。证明对任意$a,b$，必有$\eta(ab)=\eta(a)\eta(b)$或$\eta(ab)=\eta(b)\eta(a)$。
\end{problem}
\begin{proof}[补充证明]
记$A=\eta(a)$、$B=\eta(b)$、$C=\eta(ab)$。在习题26的三重乘积极化式中，依次取外侧两个变量为$ab,a$，中间变量为$b$，得到
\[
\eta(ab^2a+abab)=CBA+ABC.
\]
另一方面，$\eta(ab^2a)=AB^2A$，且$\eta(abab)=\eta((ab)^2)=C^2$。故
\[
C^2+AB^2A=CBA+ABC.
\]
移项并因式分解，
\[
(C-AB)(C-BA)=0.
\]
由于$R'$无零因子，$C=AB$或$C=BA$。
\end{proof}

\begin{problem}
\textbf{习题28：}设$\eta:R\to R'$是除环之间的加法映射，$\eta(1)=1'$，且对每对$a,b$都有$\eta(ab)=\eta(a)\eta(b)$或$\eta(b)\eta(a)$。证明$\eta$为同态或反同态。
\end{problem}
\begin{proof}[补充证明]
固定$b\in R$，记
\begin{align*}
H_b&=\{a\in R:\eta(ab)=\eta(a)\eta(b)\},\\
K_b&=\{a\in R:\eta(ab)=\eta(b)\eta(a)\}.
\end{align*}
由$\eta$的加法性，$H_b,K_b$均为$R$加法群的子群。题设说明$R=H_b\cup K_b$。一个群不能是两个真子群的并，故对每个$b$，要么$H_b=R$，要么$K_b=R$。

再令$B_+=\{b:H_b=R\}$、$B_-=\{b:K_b=R\}$。它们也是$R$加法群的子群，且$R=B_+\cup B_-$。同理，$B_+=R$或$B_-=R$。前者说明$\eta$是同态，后者说明它是反同态。
\end{proof}

\begin{problem}
\textbf{习题29：}证明除环到无零因子环的若尔当同态必为同态或反同态。
\end{problem}
\begin{proof}[补充证明]
由习题27，对每一对$a,b$，$\eta(ab)$等于$\eta(a)\eta(b)$或$\eta(b)\eta(a)$。习题28的两个加法子群论证只用到$\eta$的加法性和这逐对的二择一性质，不要求目标环为除环。因此它也适用于此处，并给出$\eta$为同态或反同态。
\end{proof}

\begin{problem}
\textbf{习题30：}设$R=\mathbb Z/(p^n)$，其中$p$为素数、$n>1$。确定零因子、素理想及相应商域。
\end{problem}
\begin{proof}[补充解答]
整数$a$在$R$中可逆当且仅当$p\nmid a$。若$p\mid a$，则$a$乘以某个非零的$p$幂即为$0$，故非单位恰为零因子。理想由$p^n$的因子对应，唯一的素理想为$P=(p)/(p^n)$。约化模$p$给出
\[
R/P\cong\mathbb Z/(p),
\]
它是域。
\end{proof}

\begin{problem}
\textbf{习题31：}设$R=\mathbb Z/(p^n)$，确定$R[x]$的单位和零因子。
\end{problem}
\begin{proof}[补充解答]
写$f=a_0+a_1x+\cdots+a_mx^m$。由习题43和45，$f$可逆当且仅当$a_0$不被$p$整除、$a_i$（$i>0$）均被$p$整除。

若非零$f$的所有系数都被$p$整除，则$p^{n-1}f=0$，故$f$是零因子。反之，设$a_j$是$f$的最低次单位系数。任取非零$g=\sum b_kx^k$，令$s$为$b_k$的最小$p$进赋值，并取具有此赋值的最低次系数$b_k$。在$fg$的$x^{j+k}$项中，$a_jb_k$的赋值恰为$s$；其余乘积或使用次数小于$j$的$f$系数，或使用次数小于$k$的$g$系数，赋值都大于$s$。该项因而非零，$fg\ne0$。所以非零零因子恰为全部系数均被$p$整除的多项式。
\end{proof}

\begin{problem}
\textbf{习题32：}设$R=R_1\oplus\cdots\oplus R_r$，$N\lhd R$。证明$N=\bigoplus_i(N\cap R_i)$，并求$\mathbb Z/(n)$的最大幂零理想和单位群。
\end{problem}
\begin{proof}[补充证明]
设$e_i$是第$i$个直和因子的单位元。任取$x\in N$，
\[
x=\sum_i e_ix,
\]
其中$e_ix\in N\cap R_i$，且分解唯一。因此得到所述直和。

若$n=\prod_i p_i^{a_i}$，中国剩余定理给出$\mathbb Z/(n)\cong\prod_i\mathbb Z/(p_i^{a_i})$。每一因子的幂零元恰为$p_i$的倍数，故最大幂零理想为$(\prod_i p_i)/(n)$；单位群为
\[
(\mathbb Z/(n))^\times\cong\prod_i(\mathbb Z/(p_i^{a_i}))^\times,
\qquad |(\mathbb Z/(n))^\times|=n\prod_{p\mid n}\left(1-\frac1p\right).
\]
\end{proof}

\subsection*{多项式与有限域（习题33--51）}

\begin{problem}
\textbf{习题33：}证明$x^3-x$在$\mathbb Z/(6)$中有六个根。
\end{problem}
\begin{proof}[补充证明]
对任意整数$a$，三个连续整数$a-1,a,a+1$的乘积同时被$2$和$3$整除。因此$6\mid a^3-a$。$\mathbb Z/(6)$的全部六个元素都是根。
\end{proof}

\begin{problem}
\textbf{习题34：}证明$x^2+1$在四元数体$H$中有无穷多个根。
\end{problem}
\begin{proof}[补充证明]
若$q=bI+cJ+dK$且$b^2+c^2+d^2=1$，四元数乘法关系使交叉项成对抵消，故
\[
q^2=-(b^2+c^2+d^2)=-1.
\]
单位球面上的三元组有无穷多个，因此根有无穷多个。
\end{proof}

\begin{problem}
\textbf{习题35：}设$u=\sqrt2+\sqrt3$，求其在$\mathbb Q$上的极小多项式。
\end{problem}
\begin{proof}[补充解答]
由$u^2=5+2\sqrt6$，得
\[
(u^2-5)^2=24,
\qquad u^4-10u^2+1=0.
\]
又$u^{-1}=\sqrt3-\sqrt2$，故$\sqrt3=(u+u^{-1})/2$、$\sqrt2=(u-u^{-1})/2$均属于$\mathbb Q(u)$。于是$\mathbb Q(u)=\mathbb Q(\sqrt2,\sqrt3)$，其次数为$4$。故极小多项式为$x^4-10x^2+1$。
\end{proof}

\begin{problem}
\textbf{习题36：}设$u^3-u+1=0$。将$(5u^2+3u-1)(2u^2-2u+6)$及$(3u^2-u+2)^{-1}$写成次数不超过$2$的多项式。
\end{problem}
\begin{proof}[补充解答]
用$u^3=u-1$、$u^4=u^2-u$逐项约去高次幂，得到
\[
(5u^2+3u-1)(2u^2-2u+6)=32u^2+6u-2.
\]
再直接相乘并约去高次幂，可验证
\[
(3u^2-u+2)(-2u^2-u+3)=7.
\]
因此$(3u^2-u+2)^{-1}=(-2u^2-u+3)/7$。
\end{proof}

\begin{problem}
\textbf{习题37：}证明$\mathbb F_2[x]/(x^3+x^2+1)$是八元域。
\end{problem}
\begin{proof}[补充证明]
三次多项式在域上不可约当且仅当它没有根。代入$0,1$分别得到$1,1$，所以$x^3+x^2+1$不可约。商环是域，每个剩余类唯一写作$a+bx+cx^2$，其中$a,b,c\in\mathbb F_2$，故有$2^3=8$个元素。
\end{proof}

\begin{problem}
\textbf{习题38：}列出$\mathbb F_5[x]$中全部首项系数为$1$的二次不可约多项式。
\end{problem}
\begin{proof}[补充解答]
$x^2+bx+c$不可约当且仅当判别式$b^2-4c$不是$\mathbb F_5$中的平方，即等于$2$或$3$。逐一代入$b=0,1,2,3,4$，得到
\begin{align*}
&x^2+2,\ x^2+3,\ x^2+x+1,\ x^2+x+2,\ x^2+2x+3,\\
&x^2+2x+4,\ x^2+3x+3,\ x^2+3x+4,\ x^2+4x+1,\ x^2+4x+2.
\end{align*}
\end{proof}

\begin{problem}
\textbf{习题39：}构造一个有$25$个元素的域。
\end{problem}
\begin{proof}[补充解答]
习题38表明$x^2+2$在$\mathbb F_5[x]$中不可约。故
\[
F=\mathbb F_5[x]/(x^2+2)
\]
是域，其元素唯一写作$a+b\bar x$（$a,b\in\mathbb F_5$），共有$5^2=25$个。
\end{proof}

\begin{problem}
\textbf{习题40：}证明$(3,x^3+2x^2+2x-1)$是$\mathbb Z[x]$的极大理想。
\end{problem}
\begin{proof}[补充证明]
模$3$约化后，多项式成为$h=x^3-x^2-x-1\in\mathbb F_3[x]$。代入$0,1,2$分别得$2,1,1$，所以三次多项式$h$无根，因而不可约。于是
\[
\mathbb Z[x]/(3,x^3+2x^2+2x-1)\cong\mathbb F_3[x]/(h)
\]
是域，故该理想极大。
\end{proof}

\begin{problem}
\textbf{习题41：}令$I=(3,x^3+2x^2+2x-1)$。证明$\mathbb Z[x]/I\cong\mathbb F_3[x]/(x^3-x^2-x-1)$，并推出$I$极大。
\end{problem}
\begin{proof}[补充证明]
系数模$3$的满同态$\rho:\mathbb Z[x]\to\mathbb F_3[x]$以$(3)$为核，并把$x^3+2x^2+2x-1$送到$x^3-x^2-x-1$。同态定理给出所述同构。由习题40中该三次式的不可约性，右侧是域，故$I$极大。
\end{proof}

\begin{problem}
\textbf{习题42：}设$F$为域且$\deg f>0$。证明$F[x]/(f)$无零因子当且仅当$f$不可约。
\end{problem}
\begin{proof}[补充证明]
若$f$不可约，则$(f)$为极大理想，商环为域。反之，若$f=gh$且$0<\deg g,\deg h<\deg f$，则$\bar g,\bar h$在商环中均非零，但$\bar g\bar h=0$，矛盾。
\end{proof}

\begin{problem}
\textbf{习题43：}设$R$为交换环。若$f\in R[x]$幂零，证明$f$的每个系数均幂零。
\end{problem}
\begin{proof}[补充证明]
写$f=a_0+a_1x+\cdots+a_mx^m$，设$f^N=0$。最高项系数为$a_m^N$，故$a_m$幂零。模去幂零理想$(a_m)$后，$f$的次数下降，且仍幂零。对次数归纳，所有$a_i$在$R/(a_m)$中幂零；于是某个$a_i^k\in(a_m)$，而$(a_m)$为幂零理想，故$a_i$也幂零。
\end{proof}

\begin{problem}
\textbf{习题44：}设$R$为交换环。证明$f\in R[x]$为零因子当且仅当存在$0\ne a\in R$使$af=0$。
\end{problem}
\begin{proof}[补充证明]
充分性直接成立。若$fg=0$且$g\ne0$，在所有被$f$消去的非零多项式中选次数最低者$g$。设$f=a_0+\cdots+a_mx^m$。因$fg=0$，最高次项系数给出$a_m$乘以$g$的最高次系数为零，故$\deg(a_mg)<\deg g$。同时$f(a_mg)=a_m(fg)=0$，由$g$的最小性得$a_mg=0$。于是$(f-a_mx^m)g=0$。依次降次，得到$a_ig=0$对每个$i$成立。取$g$的任一非零系数$a$，便有$af=0$。
\end{proof}

\begin{problem}
\textbf{习题45：}设$R$为交换环，$f=a_0x^n+a_1x^{n-1}+\cdots+a_n$。若$a_n$是单位而$a_0,\ldots,a_{n-1}$均幂零，证明$f$是$R[x]$的单位。
\end{problem}
\begin{proof}[补充证明]
令$h=a_n^{-1}(f-a_n)$。它只有有限个非零系数，且每个系数都幂零；在足够高的幂$h^N$的任一项中，必有某个系数出现到其幂零次数，故$h^N=0$。于是
\[
f^{-1}=a_n^{-1}(1+h)^{-1}
=a_n^{-1}\sum_{j=0}^{N-1}(-h)^j.
\]
\end{proof}

\begin{problem}
\textbf{习题46：}证明习题43的逆命题。
\end{problem}
\begin{proof}[补充证明]
若$f=\sum_{i=0}^m a_ix^i$且$a_i^{n_i}=0$，取$N=1+\sum_{i=0}^m(n_i-1)$。$f^N$的每一项都含$N$个系数因子；由抽屉原理，其中某个$a_i$至少出现$n_i$次，所以该项为零。故$f^N=0$。
\end{proof}

\begin{problem}
\textbf{习题47：}设$R'=R[x_1,\ldots,x_n]$，$r(R)$为$R$的幂零根。证明$r(R')=r(R)[x_1,\ldots,x_n]$。
\end{problem}
\begin{proof}[补充证明]
对变量个数归纳。习题43和46说明一元多项式幂零当且仅当各系数幂零，即$r(R[x])=r(R)[x]$。将$R[x_1,\ldots,x_{n-1}]$代替$R$，重复这一等式，便得结论。
\end{proof}

\begin{problem}
\textbf{习题48：}证明有限环中右逆也是左逆，并判断无限环是否仍成立。
\end{problem}
\begin{proof}[补充证明]
若$uv=1$且$R$有限，映射$x\mapsto ux$有右逆映射$x\mapsto vx$，故是满射，因有限而单射。由$u(vu-1)=0$得$vu=1$。无限环不成立：在可数维向量空间的自同态环中，左移算子$u$与右移算子$v$满足$uv=1$而$vu\ne1$。
\end{proof}

\begin{problem}
\textbf{习题49：}证明交换环$R$的多项式环$R[x]$中，强拟正则元恰为幂零元。
\end{problem}
\begin{proof}[补充证明]
若$f$强拟正则，则$1-xf$是$R[x]$的单位。多项式单位的正次数系数均幂零，故$f$的各系数均幂零。由习题46，$f$幂零。

反之，若$f^N=0$，则对任意$g\in R[x]$，$(fg)^N=0$，且
\[
(1-fg)^{-1}=\sum_{j=0}^{N-1}(fg)^j.
\]
由习题21的刻画，$f$强拟正则。
\end{proof}

\begin{problem}
\textbf{习题50：}设交换环$R$中每个$a$都有某个$n>1$使$a^n=a$。证明每个素理想均极大。
\end{problem}
\begin{proof}[补充证明]
若$P$为素理想，$R/P$是整环，且其中每个元素仍满足某个$\bar a^n=\bar a$。对$\bar a\ne0$，消去$\bar a$得$\bar a^{n-1}=1$，故$\bar a$可逆。于是$R/P$为域，$P$极大。
\end{proof}

\begin{problem}
\textbf{习题51：}设$R$为布尔环，即每个$a$满足$a^2=a$。证明$2a=0$、每个素理想极大且商域为$\mathbb F_2$、每个有限生成理想均为主理想。
\end{problem}
\begin{proof}[补充证明]
由$(a+a)^2=a+a$及$a^2=a$，展开并约去$2a$，得$2a=0$。又$(a+b)^2=a+b$给出$ab=ba$，故$R$交换。若$P$素，则$R/P$是布尔整环；每个非零元由$a^2=a$得$a=1$，所以$R/P\cong\mathbb F_2$，$P$极大。

对$a,b\in R$，记$c=a+b+ab$。因$ac=a$、$bc=b$，有$(a,b)\subseteq(c)$；又$c\in(a,b)$，故$(a,b)=(c)$。对生成元个数归纳，任一有限生成理想均为主理想。
\end{proof}

\subsection*{幂零根、有限域与形式幂级数（习题52--68）}

\begin{problem}
\textbf{习题52：}设$r(R)$为交换环$R$的幂零根。证明以下三项等价：$R$只有一个素理想；$R$的每个元素不是单位就是幂零元；$R/r(R)$是域。
\end{problem}
\begin{proof}[补充证明]
幂零根是全部素理想之交。若素理想只有$P$，则$r(R)=P$；任一极大理想均为素理想，故$P$是唯一极大理想。$P$外元素均为单位，$P$内元素均幂零。

若每个非单位均幂零，则$r(R)$恰为非单位集，商环中每个非零元都可逆，故$R/r(R)$是域。若商环是域，则$r(R)$极大；任一素理想都含它，只能等于它，故素理想唯一。
\end{proof}

\begin{problem}
\textbf{习题53：}证明只有有限多个理想的整环是域。
\end{problem}
\begin{proof}[补充证明]
对$0\ne a\in R$，链$(a)\supseteq(a^2)\supseteq\cdots$必有相邻两项相等。于是$a^m=a^{m+1}b$，在整环中消去$a^m$得$1=ab$。所以每个非零元可逆。
\end{proof}

\begin{problem}
\textbf{习题54：}设整环$R$的任意非空理想集合按包含关系都有极小元。证明$R$是域。
\end{problem}
\begin{proof}[补充证明]
对$0\ne a\in R$，在$\{(a^m):m\geqslant1\}$中取极小元$(a^m)$。由$(a^{m+1})\subseteq(a^m)$及极小性，二者相等。按习题53消去，$a$可逆。
\end{proof}

\begin{problem}
\textbf{习题55：}设$a,b,ab-1$均为单位。证明$a-b^{-1}$和$(a-b^{-1})^{-1}-a^{-1}$均为单位，并验证华罗庚恒等式。
\end{problem}
\begin{proof}[补充证明]
令$c=a-b^{-1}=(ab-1)b^{-1}$，故$c$可逆。又
\[
c^{-1}-a^{-1}=c^{-1}(a-c)a^{-1}=c^{-1}b^{-1}a^{-1}
\]
是单位。取逆并代入$c$，得到
\[
\bigl((a-b^{-1})^{-1}-a^{-1}\bigr)^{-1}=abc=aba-a.
\]
\end{proof}

\begin{problem}
\textbf{习题56：}设$D$为除环，$C$为其中心，$R$为子体且对每个$0\ne d\in D$有$dRd^{-1}\subseteq R$。证明$R=D$或$R\subseteq C$。
\end{problem}
\begin{proof}[补充证明]
若$R\ne D$，取$x\in D\setminus R$及$a\in R$。记$b=xax^{-1}$、$c=(x+1)a(x+1)^{-1}$，二者均在$R$。比较$xa=bx$与$xa+a=cx+c$，得$(b-c)x=c-a$。若$b\ne c$，则$x=(b-c)^{-1}(c-a)\in R$，矛盾。故$b=c=a$，即$x$与$a$交换。

任取$y\in R$，$x+y\notin R$，故$x+y$也与$a$交换，从而$y$与$a$交换。于是$a$与$D$的每个元素交换，即$a\in C$。故$R\subseteq C$。
\end{proof}

\begin{problem}
\textbf{习题57：}设$R$为实系数多项式函数环。证明每个$t\in\mathbb R$的代入映射$\eta_t(f)=f(t)$为到$\mathbb R$的同态，且每个这样的同态均为某个$\eta_t$。
\end{problem}
\begin{proof}[补充证明]
代入保持$1$、加法及乘法。设$\eta:R\to\mathbb R$为保幺同态。对实常数$c>0$，写$c=d^2$，则$\eta(c)=\eta(d)^2\geqslant0$；故$\eta$在常数域上保序。它固定有理数，由有理数的稠密性，$\eta(c)=c$对每个实数$c$成立。令$t=\eta(x)$，则
\[
\eta\left(\sum_i a_ix^i\right)=\sum_i a_it^i.
\]
\end{proof}

\begin{problem}
\textbf{习题58：}设$F$为无限域。若$g\ne0$，且$f$在所有满足$g(c_1,\ldots,c_n)\ne0$的点上为零，证明$f=0$。
\end{problem}
\begin{proof}[补充证明]
在$g$非零的点上$fg=0$，在$g$为零的点上亦然。故$fg$在$F^n$上恒为零。无限域上的多项式若处处为零，则为零多项式；由多项式环无零因子且$g\ne0$，得$f=0$。
\end{proof}

\begin{problem}
\textbf{习题59：}设$F$有$q$个元素。证明$g_0=\prod_{i=1}^n(1-x_i^{q-1})$在原点取值$1$，在其余点取值$0$。
\end{problem}
\begin{proof}[补充证明]
在原点，各因子均为$1$。若$a=(a_1,\ldots,a_n)\ne0$，某个$a_i\ne0$；由有限域乘法群的阶为$q-1$，$a_i^{q-1}=1$，相应因子为$0$。
\end{proof}

\begin{problem}
\textbf{习题60：}若$f$在原点为零而在其余点均非零，证明$g=1-f^{q-1}$与习题59的$g_0$取值相同，并推出$\deg f\geqslant n$。
\end{problem}
\begin{proof}[补充证明]
由$z^{q-1}=1$对每个$z\in F^\times$成立，$g$在原点为$1$，在其余点为$0$。将多项式逐变量模$x_i^q-x_i$约化，不增加总次数，可得到各变量次数小于$q$的唯一函数代表。$g_0$就是此代表，且其单项式$x_1^{q-1}\cdots x_n^{q-1}$系数为$(-1)^n\ne0$。因此$\deg g\geqslant n(q-1)$；又$\deg g\leqslant(q-1)\deg f$，故$\deg f\geqslant n$。
\end{proof}

\begin{problem}
\textbf{习题61：}设$F$是$q$元域，$\deg f<n$且$f(0)=0$。证明$f$还有一个非零零点。
\end{problem}
\begin{proof}[补充证明]
记$N$为$f$在$F^n$中的零点数。在$F$中计算
\[
N=\sum_{a\in F^n}\bigl(1-f(a)^{q-1}\bigr).
\]
对$k\geqslant0$，$\sum_{t\in F}t^k$只有在$k$为$q-1$的正倍数时才可能非零；当$k=0$时该和为$q=0$。故一个单项式对全体点求和后可能非零，要求每个变量的指数至少为$q-1$，总次数至少$n(q-1)$。但$\deg f^{q-1}<n(q-1)$，故$N=0$在$F$中，即整数$N$被特征$p$整除。原点已是零点，故$N\geqslant p\geqslant2$。
\end{proof}

\begin{problem}
\textbf{习题62：}设$R$为交换环、$f=\sum_{i\geqslant0}a_ix^i\in R[[x]]$。证明：(i) $f$可逆当且仅当$a_0$可逆；(ii) $R[[x]]$为整环当且仅当$R$为整环；(iii) 若$R$为整环且$F$为其商域，则$R[[x]]$商域的非零元均可写作$x^\nu(a_0+a_1x+\cdots)$，其中$\nu\in\mathbb Z$、$a_i\in F$、$a_0\ne0$。
\end{problem}
\begin{proof}[补充证明]
设$f=\sum_{i\geqslant0}a_ix^i$。若$fg=1$，常数项给出$a_0b_0=1$。反之，若$a_0$可逆，递归取
\[
b_0=a_0^{-1},\qquad b_m=-a_0^{-1}\sum_{i=1}^m a_i b_{m-i}\quad(m\geqslant1),
\]
则$g=\sum b_ix^i$满足$fg=1$。

若$R$为整环，非零幂级数的最低非零项相乘仍非零，故$R[[x]]$为整环；反向由常数嵌入可知。设$F$为$R$的商域，$R[[x]]$的商域自然嵌入$F((x))$。因此其中每个非零元均可写作
\[
x^\nu(a_0+a_1x+a_2x^2+\cdots),\qquad \nu\in\mathbb Z,\ a_i\in F,\ a_0\ne0.
\]
当$R=F$本身是域时，$F[[x]]$的商域恰为$F((x))$。
\end{proof}

\begin{problem}
\textbf{习题63：}若交换环$R$中每个非单位$a$都使$1+a$为单位，证明全部非单位构成唯一的极大理想，并判断这种环是否存在。
\end{problem}
\begin{proof}[补充证明]
记非单位集为$M$。若$a\in M$、$r\in R$而$ra$为单位，则$a$也为单位，矛盾；故$ra\in M$。若$a,b\in M$而$a+b=u$为单位，则$a/u,b/u$仍非单位。由题设，$1-a/u=b/u$应为单位，矛盾。因此$M$对加法及取负封闭，是理想。它包含所有真理想，因而是唯一极大理想。例如$\mathbb Z/(p^n)$满足该性质。
\end{proof}

\begin{problem}
\textbf{习题64：}设交换整环$R$中每个非零非单位$a$都使$1+a$仍非单位。证明全部单位加上$0$构成子域$F$，且$R\setminus F$中的元素都超越于$F$。
\end{problem}
\begin{proof}[补充证明]
记$F=R^\times\cup\{0\}$。若单位$u,v$之和非零且非单位，令$a=-(u+v)u^{-1}$，则$a$非零且非单位；但$1+a=-vu^{-1}$为单位，矛盾。所以$F$对加法封闭；其余域运算显然封闭。

若$z\in R\setminus F$在$F$上代数，取常数项非零的非零关系$c_0+c_1z+\cdots+c_mz^m=0$；若原关系常数项为零，可在整环中消去适当的$z$幂。于是
\[
z(c_1+\cdots+c_mz^{m-1})=-c_0\in F^\times,
\]
说明$z$可逆，矛盾。
\end{proof}

\begin{problem}
\textbf{习题65：}设$F$为无限域。证明不存在非常数多项式$f\in F[x]$满足$f(x)f(y)=f(x+y)$。
\end{problem}
\begin{proof}[补充证明]
设$\deg f=d>0$、首项系数为$c\ne0$。等式左边含非零项$c^2x^dy^d$，总次数为$2d$；右边总次数为$d$，矛盾。
\end{proof}

\begin{problem}
\textbf{习题66：}设$F$为域，证明$M_n(F)$的每个左理想都是主理想。
\end{problem}
\begin{proof}[补充证明]
设$L$为左理想，把$L$中各矩阵的行向量组成的子空间记为$W\subseteq F^n$。左乘矩阵单位可把任一行放入任意指定的行，其余行置零；故$L$恰是所有行向量均属于$W$的矩阵。取矩阵$A$，使其非零行构成$W$的一组基，则$L=M_n(F)A$。
\end{proof}

\begin{problem}
\textbf{习题67：}设$G$为阿贝尔群，证明$\operatorname{End}(G)$以逐点加法和映射复合构成含幺环。
\end{problem}
\begin{proof}[补充证明]
若$\eta,\xi$为自同态，则
\[
(\eta+\xi)(a+b)=\eta(a)+\eta(b)+\xi(a)+\xi(b)
=(\eta+\xi)(a)+(\eta+\xi)(b),
\]
重排使用$G$交换。零映射、逐点负映射给出加法零元与逆元；加法结合且交换。复合结合，恒等映射为乘法单位。自同态的加法性给出复合对逐点加法的左右分配律。
\end{proof}

\begin{problem}
\textbf{习题68：}分别求$\mathbb Z^+$、无限循环群、$\mathbb Q^+$及$n$个$p$阶循环群的直和的自同态环。
\end{problem}
\begin{proof}[补充解答]
$\mathbb Z^+$及任一无限循环群的自同态由生成元的像唯一决定；该像是生成元的整数倍，故自同态环均同构于$\mathbb Z$。

$\mathbb Q^+$的自同态$\eta$由$\eta(1)=r\in\mathbb Q$唯一决定，因为$n\eta(m/n)=m\eta(1)$。故$\eta(q)=rq$，自同态环同构于$\mathbb Q$。

$n$个$p$阶循环群的直和是$\mathbb F_p$上的$n$维向量空间，自同态恰为线性变换，所以其自同态环同构于$M_n(\mathbb F_p)$。
\end{proof}


%
